\documentclass[11pt,a4paper]{amsart}
\usepackage[utf8]{inputenc}
\usepackage[T1]{fontenc}
\usepackage[french]{babel}
\usepackage{amsmath,amssymb,amsthm}
\usepackage{mathrsfs}
\usepackage{graphicx}
\usepackage[all]{xy}
\usepackage{array}
\usepackage{enumitem}

% Theorem environments
\theoremstyle{plain}
\newtheorem{theorem}{Th\'eor\`eme}[section]
\newtheorem{lemma}[theorem]{Lemme}
\newtheorem{proposition}[theorem]{Proposition}
\newtheorem{corollary}[theorem]{Corollaire}

\theoremstyle{definition}
\newtheorem{definition}[theorem]{D\'efinition}
\newtheorem{remark}[theorem]{Remarque}
\newtheorem{remarks}[theorem]{Remarques}
\newtheorem{construction}[theorem]{Construction}
\newtheorem{example}[theorem]{Exemple}
\newtheorem{variante}[theorem]{Variante}

\DeclareMathOperator{\Pic}{Pic}
\DeclareMathOperator{\Aut}{Aut}
\DeclareMathOperator{\Isom}{Isom}
\DeclareMathOperator{\Spec}{Spec}
\DeclareMathOperator{\Proj}{Proj}
\DeclareMathOperator{\Ker}{Ker}
\DeclareMathOperator{\reg}{reg}
\DeclareMathOperator{\supp}{supp}
\DeclareMathOperator{\Sch}{Sch}
\DeclareMathOperator{\Gl}{GL}
\DeclareMathOperator{\SL}{SL}
\DeclareMathOperator{\Ext}{Ext}
\DeclareMathOperator{\Hom}{Hom}
\DeclareMathOperator{\End}{End}
\DeclareMathOperator{\Lie}{Lie}
\DeclareMathOperator{\Def}{Def}

\newcommand{\CC}{\mathbb{C}}
\newcommand{\CH}{\mathcal{H}}
\newcommand{\CO}{\mathcal{O}}
\newcommand{\CG}{\mathcal{G}}
\newcommand{\CF}{\mathcal{F}}
\newcommand{\CM}{\mathcal{M}}
\newcommand{\CT}{\mathcal{T}}
\newcommand{\CL}{\mathcal{L}}
\newcommand{\CS}{\mathcal{S}}
\newcommand{\CA}{\mathcal{A}}
\newcommand{\CI}{\mathcal{I}}
\newcommand{\CJ}{\mathcal{J}}
\newcommand{\CU}{\mathcal{U}}
\newcommand{\CV}{\mathcal{V}}
\newcommand{\mll}{\mathcal{M}}
\DeclareMathOperator{\Id}{Id}
\DeclareMathOperator{\Ob}{Ob}
\DeclareMathOperator{\im}{im}
\DeclareMathOperator{\car}{car}
\newcommand{\uHom}{\underline{\Hom}}

\newcommand{\Gm}{\mathbb{G}_m}
\newcommand{\ZZ}{\mathbb{Z}}
\newcommand{\QQ}{\mathbb{Q}}
\newcommand{\CCb}{\mathbb{C}}
\newcommand{\FF}{\mathbb{F}}
\newcommand{\PP}{\mathbb{P}}
\newcommand{\HH}{\mathbb{H}}
\newcommand{\RR}{\mathbb{R}}
\newcommand{\Af}{\mathbb{A}}

\newcommand{\fppf}{\textup{fppf}}
\newcommand{\fpqc}{\textup{fpqc}}
\newcommand{\et}{\textup{\'et}}
\newcommand{\ens}{\textup{Ens}}
\newcommand{\mg}{\mathfrak{m}}

\begin{document}

\section*{Chapitre II. Courbes elliptiques g\'en\'eralis\'ees (suite)}

\setcounter{section}{3}
\setcounter{subsection}{1}

\subsection{}

\noindent\textbf{(3.1.1)} $\hspace{3em} G$ agit trivialement sur $\Pic^0(C/S)$.

\bigskip

\noindent\textbf{(3.1.2)} $\hspace{3em}$ Pour chaque point g\'eom\'etrique $s$ de $S$, $G(\bar{s})$ agit transitivement sur l'ensemble des composantes irr\'eductibles de $C_{\bar{s}}$.

\medskip

Dans le cas particulier important o\`u $G$ est le sch\'ema en groupes constant $\ZZ$, l'action de $G$ est d\'efinie par l'automorphisme $g = \tau(1)$ de $C$ et (3.1.1), (3.1.2) deviennent

\medskip

\noindent\textbf{(3.1.3)} $\hspace{3em} g$ agit trivialement sur $\Pic^0(C/S)$.

\bigskip

\noindent\textbf{(3.1.4)} $\hspace{3em}$ Pour chaque point g\'eom\'etrique $s$ de $S$, $g$ permute transitivement les composantes irr\'eductibles de $C_{\bar{s}}$.

\begin{theorem}[3.2]
Soit $G$ agissant sur $C$ comme en 3.1. Si les conditions (3.1.1) et (3.1.2) sont v\'erifi\'ees, il existe sur $C$ une et une seule structure de courbe elliptique g\'en\'eralis\'ee d'unit\'e $e$ telle que $G$ agisse sur $C$ par translations. De plus, tout automorphisme de $C$ sur $S$ qui commute \`a $G$ et v\'erifie (3.1.1) est une translation.
\end{theorem}

\noindent\textit{Plan de la d\'emonstration.} Nous prouverons les \'enonc\'es suivants.

\medskip

\noindent\textbf{(A)} Soit $s$ un point g\'eom\'etrique de $S$ tel que $G_{\bar{s}}$ agissant sur $C_{\bar{s}}$ v\'erifie (3.1.1) et (3.1.2). Il existe un voisinage \fppf\ $U$ de $s$ et $g \in G(U)$ tel que l'action de $g$ sur $C_U$ v\'erifie (3.1.3), (3.1.4).

\medskip

\noindent\textbf{(B)} 3.2 pour $G = \ZZ$.

\medskip

Montrons d\'ej\`a que (A) et (B) impliquent 3.2. Les probl\`emes sont locaux pour la topologie \fppf. Si $G$ agissant sur $C_{\bar{s}}$ v\'erifie (3.1.1), (3.1.2), on peut donc supposer qu'il existe $g \in G(S)$ comme en (A). D'apr\`es (B), il existe une et une seule structure de courbe elliptique g\'en\'eralis\'ee d'unit\'e $e$ telle que $g$ soit une translation. De plus, puisque $G$ et $g$ commutent et que $G$ v\'erifie (3.1.1), $G$ agit par translations. Enfin, tout automorphisme v\'erifiant (3.1.1) qui commute \`a $G$ commute \`a $g$, donc est une translation.

\medskip

\noindent\textit{Preuve de (A).} Si $C_{\bar{s}}$ est lisse, on prend pour $U$ l'ouvert de $S$ au-dessus duquel $C$ est lisse, et $g = e$. Sinon, $C_{\bar{s}}$ est un $n$-gone et l'image de $G(\bar{s})$ dans le groupe des automorphismes du diagramme $\Gamma(C_{\bar{s}})$ des composantes irr\'eductibles de $C_{\bar{s}}$ est d'ordre $\leq n$. D'apr\`es (3.1.1), $G(\bar{s})$ agit par rotations sur $\Gamma(C_{\bar{s}})$, le groupe des rotations de $\Gamma$ \'etant cyclique d'ordre $n$, il existe $g_{\bar{s}} \in G(\bar{s})$ d'image un g\'en\'erateur de ce groupe, et $g_{\bar{s}}$ permute transitivement les composantes irr\'eductibles de $C_{\bar{s}}$. Localement \fppf\ il existe une section $g$ de $G$ qui induise $g_{\bar{s}}$. On peut appliquer II.1.16 \`a $g$ pour v\'erifier (3.1.4) dans un voisinage de $s$. Que (3.1.3) soit v\'erifi\'e par $g$ dans un voisinage de $s$ r\'esulte de II.1.13.

\medskip

\noindent\textit{Preuve de 3.2 pour $G = \ZZ$.} Soient $\tau$ l'action de $G$ sur $C$ et $g = \tau(1)$. En appliquant II.1.16 \`a $g$, on se ram\`ene \`a supposer que les fibres g\'eom\'etriques de $C$ sont toutes lisses ou des $n$-gones ($n$ fix\'e). Dans ce cas, $g^n \cdot e$ est fibre par fibre dans la m\^eme composante connexe de $C$ que $e$ et $g^n$ est justiciable du lemme suivant.

\begin{lemma}[3.3]
Soit $h$ un automorphisme de $C$ qui commute \`a $g$ et v\'erifie (3.1.3). Si $h \cdot e$ est fibre par fibre dans la m\^eme composante irr\'eductible de $C$ que $e$, alors $h$ est de la forme $x \mapsto \lambda + x$ pour $\lambda$ une section de $\Pic^0(C/S)$.
\end{lemma}

\noindent\textit{Preuve.} On aura $\lambda = \text{classe de } \CO(h \cdot e - e)$. L'hypoth\`ese (3.1.3) assure que $x \mapsto \lambda + x$ commute \`a $g$. Rempla\c{c}ant $h$ par $h(x - \lambda)$, on se ram\`ene \`a supposer que $h \cdot e = e$. Puisque $h$ v\'erifie (3.1.3), pour toute section $\xi$ de $C^{\reg}$, on a $h(g^k(\xi) \cdot e) = g^k(\xi \cdot e)$. Toute section de $C^{\reg}$ \'etant localement pour la topologie \fppf\ de la forme $g^k(\xi)$, $h$ agit trivialement sur $C^{\reg}$, donc sur $C$.

\medskip

\noindent\textit{Preuve de 3.2 (suite).} On a $g^n \cdot x = \lambda + x$. Soit $\mu$ tel que $\lambda = \mu - n$ et posons $g_1 = g - \mu$. Alors $g_1^n = \Id$ et $g_1$ d\'efinit une action de $\ZZ/n$ sur $C$. Cette action est libre et les fibres g\'eom\'etriques de $C/(\ZZ/n)$ sont int\`egres. Sur $C/(\ZZ/n)$, il existe donc une et une seule structure de courbe elliptique g\'en\'eralis\'ee de section neutre $e$ (II.2.7). Sur $C$, il existe une et une seule structure de courbe elliptique g\'en\'eralis\'ee d'unit\'e $e$ telle que $g_1$ soit une translation (cf.~II.1.17). D'apr\`es 3.3, pour $\lambda$ une section de $\Pic^0$ et pour cette structure de courbe elliptique g\'en\'eralis\'ee, $(\lambda + e) + x = \lambda + x$. La loi obtenue est donc aussi la seule telle que $g$ soit une translation. Enfin, si $h$ commute \`a $g$ et v\'erifie (3.1.3), il existe localement $k$ tel que $hg^k$ soit justiciable de 3.3, et $h$ est une translation.

\begin{remark}[3.4]
Si $C$ est une courbe elliptique g\'en\'eralis\'ee sur $S$, on peut montrer qu'il existe un et un seul homomorphisme
\[
u : C^{\reg} \longrightarrow \Aut(C)
\]
tel que, apr\`es tout changement de base, pour toute section $t$ de $C^{\reg}$ sur $S$ on ait $u(\xi \cdot t) \cdot e = t$.
\end{remark}


\section{Chapitre III. Th\'eor\`emes de repr\'esentabilit\'e}

\setcounter{subsection}{-1}
\subsection{Introduction et notations}

Dans ce chapitre et les suivants, nous d\'efinissons et \'etudions divers champs alg\'ebriques, classifiant les courbes elliptiques g\'en\'eralis\'ees munies de diverses structures additionnelles.

\subsubsection{}
Notons $\mll$ la cat\'egorie fibr\'ee en groupo\"ides sur la cat\'egorie $(\Sch)$ des sch\'emas suivante: $\mll(S)$ est la cat\'egorie dont les objets sont les courbes elliptiques g\'en\'eralis\'ees sur $S$, et les morphismes les $S$-isomorphismes; le foncteur ``image r\'eciproque par $u \colon S \to T$'' est $E \mapsto E \times_T S$. La cat\'egorie $\mll$ est un champ sur $(\Sch)$ pour la topologie \fpqc\ (II.2.1) mais n'est pas un champ alg\'ebrique.

\subsubsection{}
Soit $\mll^*$ le sous-champ suivant de $\mll$: $\mll^*$ classifie les courbes elliptiques g\'en\'eralis\'ees $C/S$ telles que, pour tout point g\'eom\'etrique $s$ de $S$, la caract\'eristique de $k(s)$ ne divise pas le nombre de composantes irr\'eductibles de la fibre g\'eom\'etrique $C_{\bar{s}}$. Le r\'esultat principal 2.5 de ce chapitre est que $\mll^*$ est un champ alg\'ebrique. De ce th\'eor\`eme technique r\'esulteront les th\'eor\`emes de repr\'esentabilit\'e des chapitres suivants. Nous le d\'eduirons des crit\`eres g\'en\'eraux de M.~Artin. Pour pouvoir appliquer ceux-ci, le point essentiel sera de prouver la prorepr\'esentabilit\'e effective de certains foncteurs; ce sera fait au \S\ 1.

\subsubsection{}
Nous adopterons les principes de notation suivants.

\begin{enumerate}[label=\alph*)]
\item $\mll$, affect\'e d'un indice, d\'esigne le champ alg\'ebrique classifiant les courbes elliptiques g\'en\'eralis\'ees, soumises \`a certaines restrictions, ou munies de structures additionnelles, dont l'esp\`ece d\'epend de l'indice.

\item L'exposant $h$ indique le sous-champ ouvert correspondant aux courbes elliptiques proprement dites.

\item L'exposant $h$ indique le sous-champ ouvert obtenu en \^otant les points correspondant aux courbes elliptiques supersinguli\`eres de caract\'eristique $p \nmid n$, o\`u $n$ d\'epend du contexte.

\item L'exposant $\infty$ indique le sous-champ ferm\'e image du lieu singulier de la courbe universelle.

\item Une notation $\mll[1/n]$ indique soit qu'on consid\`ere un champ sur $\ZZ[1/n]$, soit, si $\mll$ a d\'ej\`a \'et\'e d\'efini, d\'esigne $\mll \otimes_{\ZZ} \ZZ[1/n]$.

\item Le sch\'ema grossier de modules d\'efini par un champ alg\'ebrique (I.1.8) se note par le m\^eme symbole, avec $\mll$ remplac\'e par $M$.
\end{enumerate}

Nous \'etendrons ces notations \`a certains champs au-dessus de $\ZZ$, non n\'ecessairement d\'efinis par un ``probl\`eme de modules'' stricto sensu.


\subsection{Un th\'eor\`eme de prorepr\'esentabilit\'e}

\subsubsection{}
Soient $\Lambda$ un anneau local complet de corps r\'esiduel $k$ et $C_0$ une courbe elliptique g\'en\'eralis\'ee sur $k$. On suppose que $C_0$ est lisse ou l'image r\'eciproque sur $k$ de la courbe elliptique g\'en\'eralis\'ee standard \`a $n$ c\^ot\'es (II.1.9). Soit $D$ le foncteur des d\'eformations de $C_0$: pour $A$ une $\Lambda$-alg\`ebre noeth\'erienne locale compl\`ete de corps r\'esiduel $k$, $D(A)$ est l'ensemble des classes d'isomorphie de courbes elliptiques g\'en\'eralis\'ees $C/\Spec(A)$, munies d'un isomorphisme $C \otimes_A k \simeq C_0$. Ce \S\ est consacr\'e \`a la d\'emonstration du th\'eor\`eme suivant.

\begin{theorem}[1.2]
Supposons $C_0$ lisse ou $n$ premier \`a l'exposant caract\'eristique $p$ de $k$. On a alors

\begin{enumerate}[label=(\roman*)]
\item $(C_0,+)$ n'a pas d'automorphismes infinit\'esimaux (non nuls).
\item Le foncteur des d\'eformations $D$ de $(C_0,+)$ est effectivement prorepr\'esentable.
\item $D$ est prorepr\'esent\'e par une courbe $C$ sur $\Spec(\Lambda[[t]])$.
\item Si $C_0$ est non lisse, on peut choisir $t$ de sorte que $t=0$ soit l'image du sous-sch\'ema de non lissit\'e.
\end{enumerate}
\end{theorem}

Les automorphismes infinit\'esimaux de $(C_0,+)$ forment le groupe $\Lie \Aut(C_0,+)$, de sorte que (i) r\'esulte de (II.1.10). Il r\'esulte de (i) et de [25] que $D$ est prorepr\'esentable.

\subsubsection{Le cas o\`u $C_0$ est irr\'eductible}

Dans ce cas, le foncteur des d\'eformations de $(C_0,+)$ s'identifie par II.2.7 au foncteur $D'$ des d\'eformations de $(C_0,e)$, le couple form\'e de la courbe $C_0$ et de la section marqu\'ee $e$. Ce dernier foncteur est effectivement prorepr\'esentable, un sch\'ema formel en courbes compl\`etes \'etant toujours projectif, donc alg\'ebrisable. Que $D'$ v\'erifie (iii) r\'esulte des deux formules suivantes:
\begin{equation}\label{eq:1.3.1}
\dim \Ext^1(\CO_{C_0}(e), \CO_{C_0}) = 1
\end{equation}
\begin{equation}\label{eq:1.3.2}
\dim \Ext^2(\CO_{C_0}(e), \CO_{C_0}) = 0
\end{equation}

En effet, le $\Ext^2$ est le groupe o\`u vivent les obstructions aux d\'eformations. S'il s'annule, $D'$ est prorepr\'esent\'e par une alg\`ebre de s\'eries formelles sur $\Lambda$. Le $\Ext^1$ s'identifie \`a $D'(k[\varepsilon]/(\varepsilon^2))$, et d\'etermine la dimension relative.

Prouvons (1.3.1) et (1.3.2). Puisque $\omega_{C_0}$ est isomorphe \`a $\CO_{C_0}$, on a par dualit\'e (I.2.2)
\[
\dim \Ext^i(\CO_{C_0}(e), \CO_{C_0}) = \dim H^{1-i}(C_0, \CO_{C_0}(e)).
\]
La formule (1.3.2) en r\'esulte trivialement. Pour (1.3.1), on peut supposer $k$ alg\'ebriquement clos.

Si $C_0$ est lisse, $\CO_{C_0}(e) \simeq \omega_{C_0}(e) \simeq \CO_{C_0}(e)$ est un faisceau inversible de degr\'e 1, d'o\`u $\dim H^0(C_0, \CO_{C_0}(e)) = 1$ dans ce cas.

Sinon, le normalis\'e $\pi \colon \tilde{C}_0 \to C_0$ de $C_0$ est isomorphe \`a $\PP^1$ et un calcul local montre que le morphisme canonique
\[
\CO_{C_0} \longrightarrow \pi_* \CO_{\tilde{C}_0}
\]
est surjectif \`a noyau de longueur un, concentr\'e au seul point singulier de $C_0$. L'assertion dans ce cas r\'esulte de la suite exacte de cohomologie: on utilise que
\[
H^0(\tilde{C}, \pi^* \CO_{C_0}(e)) \simeq H^0(\PP^1, \CO_{\PP^1}(e)) = 0 \quad (\text{car } \deg \CO_{\PP^1}(e) = -1).
\]
Cela d\'emontre (iii).

Prouvons (iv). La courbe $C_0$ est ici polygone de N\'eron \`a un c\^ot\'e. La th\'eorie des d\'eformations des points quadratiques ordinaires montre qu'il existe $u \in \Lambda[[t]]$ tel que $u=0$ soit l'image du lieu de non lissit\'e de $C$. D'apr\`es (II.1.15), sur $\Spec(\Lambda[[t]]/(u))$, $C$ est isomorphe \`a l'image r\'eciproque du polygone de N\'eron \`a un c\^ot\'e standard (II.1.1). La propri\'et\'e universelle de $\Lambda[[t]]$ implique alors que $\Spec(\Lambda[[t]]/(u))$ est non ramifi\'e sur $\Spec(\Lambda)$, donc que $\Lambda[[t]] \simeq \Lambda[[u]]$. Ceci ach\`eve la d\'emonstration pour $C_0$ irr\'eductible.

\subsubsection{Le cas g\'en\'eral}

On peut supposer que $C_0$ est le polygone de N\'eron standard \`a $n$ c\^ot\'es sur $k$, d\'eduit de $\PP^1 \times \ZZ/n$ par recollement. Soit $H_0$ le sous-groupe cyclique d'ordre $n$ de $C_0$, image de $\{1\} \times \ZZ/n$.

Soit $D''$ le foncteur des d\'eformations de $(C_0, +, H_0)$: $D''(A)$ est l'ensemble des classes d'isomorphie de syst\`emes $(C,+,H)$ ($(C,+)$ courbe elliptique g\'en\'eralis\'ee sur $A$, $H$ sous-groupe de $C^{\reg}$ isomorphe \`a $\ZZ/n$), munis d'un isomorphisme de leur fibre sp\'eciale avec $(C_0,+,H_0)$.

\begin{lemma}[1.4.1]
On a $D'' \xrightarrow{\sim} D$.
\end{lemma}

Cela r\'esulte de ce que $C_n$ est fini \'etale sur $A$ (car $n$ est premier \`a $p$).

\subsubsection{}
Dans la fin de la d\'emonstration, on ne suppose plus que $n$ est premier \`a $p$ (cette hypoth\`ese ne servira que via 1.4.1). Soit $(C,+,H)$ comme plus haut (sur $A$ local complet). Le groupe $H$ agit librement sur $C$ (car $H_0$ agit librement sur $C_0$), de sorte que $C/H$ (qui existe car $C$ est projectif) est plat sur $A$. La loi $+$ passe au quotient, et $(C/H,+)$ est une d\'eformation du polygone de N\'eron \`a un c\^ot\'e; on obtient ainsi un morphisme de $D''$ dans $D_1$.

\begin{lemma}[1.4.3]
Le morphisme 1.4.2 de $D''$ dans $D_1$ est un isomorphisme (on ne suppose pas ici $n$ premier \`a $p$).
\end{lemma}

D'apr\`es EGA IV.18.1.2 pour $A$ artinien, joint au th\'eor\`eme d'existence en g\'eom\'etrie formelle EGA III.5, pour $A$ local complet, pour toute d\'eformation $(C,+)$ de $(C_0/H_0,+)$ sur $A$, il existe une et une seule d\'eformation $C$ de la courbe $C_0$ sur $A$, munie de $\pi \colon C_0 \to C_0/H_0$. On applique alors II.1.17.

Ach\`evons la d\'emonstration de 1.2. D'apr\`es 1.4.1 et 1.4.3, $D$ est isomorphe \`a $D_1$ \'etudi\'e en 1.3 et v\'erifie donc 1.2 (ii) (iii). L'assertion (iv) r\'esulte de l'assertion analogue pour $D_1$, car l'image du sous-sch\'ema de non lissit\'e est la m\^eme pour $C$ et $C/H$ (notations de 1.4.2).


\subsection{Construction de $\mll^*$}

\begin{lemma}[2.1]
$\mll$ (d\'efini en 0.1) est un champ pour la topologie \fpqc.
\end{lemma}

\noindent\textit{Preuve:} Soit $S' \to S$ un morphisme \fpqc\ et $p' \colon C' \to S'$ une courbe elliptique g\'en\'eralis\'ee munie d'une donn\'ee de descente. Les sous-sch\'emas $S'_n$ de $S'$ d\'efinis en II.1.15 se descendent en des sous-sch\'emas $S_n$ de $S$ (descente \fpqc\ de sous-sch\'emas ferm\'es). Par localisation sur $S$, pour la topologie de Zariski, on peut donc supposer que les fibres g\'eom\'etriques de $C'$ sont lisses ou ont $n$ composantes irr\'eductibles. Le sous-sch\'ema $C'_n$ de $C'$ est alors fini sur $S'$ (cf.~II.1.20) et rencontre chaque composante irr\'eductible de chaque fibre g\'eom\'etrique de $C'$. Le faisceau inversible $\CL = \CO_{C'}(C'_n)$ est relativement ample; il est muni d'une donn\'ee de descente. L'assertion r\'esulte donc de la descente des sch\'emas quasi-projectifs (SGA 1, VIII.7.8 et SGA 1, VIII.5.2).

\medskip

Rappelons que pour $\CG$ un champ sur $(\Sch/S)$ on identifie les 1-morphismes $X \to \CG$ d'un $S$-sch\'ema $X$ dans $\CG$ aux objets de $\CG(X)$ (r\`egle $X/S$).

\subsubsection{}
M.~Artin a d\'emontr\'e un crit\`ere maniable pour prouver qu'un foncteur est repr\'esentable par un espace alg\'ebrique ([2], [3]). La m\^eme d\'emonstration fournit le crit\`ere suivant pour v\'erifier qu'un champ est alg\'ebrique.

\begin{theorem}[2.3] (M.~Artin)
Soit $S$ un sch\'ema de type fini sur un corps ou sur un anneau excellent de Dedekind. Soit $\CG$ une cat\'egorie fibr\'ee en groupo\"ides sur $(\Sch/S)$.

$\CG$ est un champ alg\'ebrique localement de type fini sur $S$ si (et seulement si)
\begin{enumerate}[label=(\arabic*)]
\item[(0)] $\CG$ est un champ pour la topologie \'etale.

\item[(1)] $\CG$ est localement de pr\'esentation finie, (i.e. pour chaque syst\`eme projectif filtrant de $S$-sch\'emas affines $\Spec(A_i)$, le foncteur canonique
\[
\varinjlim_i \CG(\Spec(A_i)) \longrightarrow \CG(\varprojlim_i \Spec(A_i))
\]
d\'efinit une \'equivalence des cat\'egories).

\item[(2)] Soient $\xi, \eta \in \Ob(\CG(X))$ deux 1-morphismes d'un $S$-sch\'ema $X$ de type fini sur $S$ vers $\CG$. Alors $\Isom_X(X; \xi, \eta)$ est un espace alg\'ebrique localement de type fini sur $S$.

\item[(3)] Soit $k_0$ un $\CS$-corps de type fini (i.e. $\Spec(k_0)$ est de type fini sur $S$), et soit $\xi_0$ un 1-morphisme de $\Spec(k_0)$ vers $\CG$. Il existe un anneau local complet $R$, un morphisme $u$ du spectre d'une extension s\'eparable finie $k'_0$ de $k_0$ dans le point ferm\'e $s$ de $\Spec(R)$, et un diagramme commutatif
\[
\xymatrix{
\Spec(k'_0) \ar[r]^{\xi_0} \ar[d] & \CG \\
\Spec(R) \ar[ru]_{\xi}
}
\]
avec $\xi$ formellement \'etale en $s$.

\item[(4)] Si $\xi$ est un 1-morphisme d'un $S$-sch\'ema de type fini $X$ dans $\CG$ et que $\xi$ est formellement \'etale en un point $x$ de type fini sur $S$, alors $\xi$ est formellement \'etale en chaque point de type fini sur $S$ dans un voisinage ouvert de $x$.
\end{enumerate}
\end{theorem}

\begin{remarks}[2.4]
\ \begin{enumerate}[label=(\roman*)]
\item Si les corps r\'esiduels de type fini de $S$ sont parfaits, pour \'etablir que $\CG$ est un champ alg\'ebrique on peut remplacer (0)--(4) par (0), (1), (2), (3$'$), (4), o\`u (3$'$) est la condition suivante:

(3$'$) Soient $s \in S$ et $k_0$ une extension finie de $k(s)$. On suppose que $u \colon \Spec(k_0) \to S$ est de type fini. Par hypoth\`ese, $k(s)$ est alors parfait et $k_0/k(s)$ s\'eparable. Il existe donc un unique anneau local complet $A(k_0)$ de corps r\'esiduel $k_0$, muni de $u \colon \Spec(A(k_0)) \to S$ formellement \'etale et prolongeant $u$.

Pour $\xi_0 \in \Ob \CG(k_0)$, nous noterons $\Def(\xi_0)$ la cat\'egorie suivante sur la duale de $\CA(k_0)$: pour $A \in \Ob \CA(k_0)$, un objet de $\Def(\xi_0)(A)$ est un objet $\xi$ de $\CG(A)$ muni d'un isomorphisme $\xi \xrightarrow{\sim} \xi_0$ (image de $\xi$ dans $\CG(k_0)$). La condition (3$'$) est la conjonction de (3$'$a) et (3$'$b):

(3$'$a) Les objets de $\CG$ n'ont pas d'automorphismes infinit\'esimaux; plus pr\'ecis\'ement, pour $k_0$ et $\xi_0$ comme ci-dessus, et $\xi'_0$ l'image r\'eciproque de $\xi_0$ sur $k_0[\varepsilon]/(\varepsilon^2)$, on a $\Aut(\xi'_0) \xrightarrow{\sim} \Aut(\xi_0)$. Cette condition, et (2), implique que les objets de $\Def(\xi_0)(A)$ n'ont pas d'automorphismes non triviaux.

(3$'$b) Pour $k_0$ et $\xi_0$ comme plus haut, le foncteur
\[
A \mapsto \text{l'ensemble des classes d'isomorphies dans } \Def(\xi_0)(A)
\]
est effectivement prorepr\'esentable (par un anneau local complet $R$ et $\xi \in \CG(R)$ se r\'eduisant selon $\xi_0$).

Il suffit m\^eme de v\'erifier (3$'$b) apr\`es extension des scalaires \`a une extension finie de $k_0$.

\item Si $S$ est de type fini sur un corps ou sur un anneau excellent de Dedekind ayant un nombre infini de points et si les anneaux locaux complets $R$ qui apparaissent dans la condition (3$'$b) sont tous normaux et de m\^eme dimension de Krull, on peut supprimer la condition (4).
\end{enumerate}

(Pour le lecteur familier avec [2], (i) est facile \`a prouver. Pour (ii), voir [2], Thm.~3.9).
\end{remarks}

\subsubsection{}
Le polygone de N\'eron \`a $n$ c\^ot\'es standard d\'efinit un morphisme
\[
f_n \colon \Spec(\ZZ) \longrightarrow \mll^*.
\]

\begin{theorem}[2.5] \ \begin{enumerate}[label=(\roman*)]
\item $\mll^*$ (d\'efini en 0.2) est un champ alg\'ebrique lisse sur $\Spec(\ZZ)$.
\item L'image $\mll^{*,\infty}$ du sch\'ema de non-lissit\'e de la courbe elliptique g\'en\'eralis\'ee universelle au-dessus de $\mll^*$ est r\'eunion des images de $f_n$ ($n \geq 1$).
\end{enumerate}
\end{theorem}

\noindent\textit{Preuve.} Pour prouver que $\mll^*$ est un champ alg\'ebrique, nous utilisons le crit\`ere 2.3 (avec $S = \Spec(\ZZ)$). Passons en revue les conditions (0) \`a (4).

(0) r\'esulte de 2.1.

(1) r\'esulte par des arguments standard de EGA IV.8.8.

(2) Soient $X$ un $S$-sch\'ema de type fini et $(p_1 \colon C_1 \to X, +)$, $(p_2 \colon C_2 \to X, +)$ deux courbes elliptiques g\'en\'eralis\'ees. Prouvons que $\Isom_X((C_1,+),(C_2,+))$ est repr\'esentable. Ce foncteur est un faisceau \fpqc\ (2.1) de sorte que le probl\`eme est local. On peut donc supposer que les fibres g\'eom\'etriques de $C_1$ sont lisses ou \`a $n$ c\^ot\'es ($n$ convenable) et qu'il existe $\alpha \colon (\ZZ/n)^2 \to C_1$ dont l'image rencontre chaque composante irr\'eductible de chaque fibre g\'eom\'etrique de $C_1$. Le morphisme $f \mapsto f \circ \alpha(1)$ envoie $\Isom$ dans le sous-sch\'ema ouvert de $\uHom_{C_1}$ des $g$ tels que les $g^k$ rencontrent toutes les composantes des fibres g\'eom\'etriques. Il r\'esulte de l'assertion suivante que ce morphisme est relativement repr\'esentable, donc $\Isom$ repr\'esentable.

\medskip

\noindent\textbf{Proposition.} \textit{Soient $C_1$, $C_2$ deux courbes elliptiques g\'en\'eralis\'ees sur $X$ \`a fibres lisses ou \`a $n$ c\^ot\'es, et munies de $\alpha_i \colon (\ZZ/n)^2 \to C_i^{\reg}$ isomorphe \`a $(\ZZ/n)^2$, rencontrant toutes les composantes de toutes les fibres g\'eom\'etriques ($i=1,2$).}

\textit{D'apr\`es II.1.17, le morphisme $\Isom_X((C_1,+),(C_2,+)) \to \Isom_X(C_1^{\reg},C_2^{\reg})$ est en effet repr\'esentable par un morphisme \'etale, tandis qu'il r\'esulte de la th\'eorie du foncteur de Hilbert que $\Isom_X((C_1/H_1,e),(C_2/H_2,e))$ est repr\'esentable.}

\medskip

(3$'$) Soit $k_0$ un $\Spec(\ZZ)$-corps de type fini; c'est donc un corps fini.

(3$'$a) r\'esulte de 1.2.(i).

(3$'$b) r\'esulte de 1.2.(ii).

(4) D'apr\`es 1.2 (iii), les foncteurs de d\'eformations sont prorepr\'esent\'es par des anneaux r\'eguliers de dimension 2. D'apr\`es 2.4 (ii), la condition (4) est automatiquement v\'erifi\'ee.

La lissit\'e de $\mll^*$ r\'esulte de sa lissit\'e formelle 1.2.(iii). La seconde assertion de 2.5 r\'esulte de II.1.15 et de 1.2.(iv).
\qed

\begin{remark}[2.6]
$\mll^*$ est horriblement non s\'epar\'e! Il r\'esultera de (IV.2.2) que le sous-champ ouvert $\mll^h_1$ de $\mll^*$ classifiant les courbes elliptiques g\'en\'eralis\'ees \`a fibres g\'eom\'etriques int\`egres est propre et lisse sur $\ZZ$. Le ``fourre-tout'' $\mll^*$ nous sera utile pour compactifier des modules de courbes elliptiques avec niveau.

Le sous-champ ouvert $\mll^h(n)[1/n]$ de $\mll^*[1/n]$ qui classifie les courbes elliptiques g\'en\'eralis\'ees \`a fibres g\'eom\'etriques lisses ou des polygones de N\'eron \`a $n$ c\^ot\'es est lui aussi propre et lisse sur $\Spec(\ZZ[1/n])$. Il est toutefois moins naturel que $\mll^h_1[1/n]$.
\end{remark}


\section{Chapitre IV. Structures de niveau}

Dans ce chapitre, nous d\'efinissons et \'etudions les sch\'emas de modules de courbes elliptiques dont les points de division par $n$ sont munis de structures additionnelles. Nous \'etudions en d\'etail la structure \`a l'infini de ces sch\'emas. Leur r\'eduction modulo $p$, pour $p \nmid n$, ne sera consid\'er\'ee s\'erieusement qu'au chapitre suivant.

\subsection{Contractions}

\subsubsection{}
Soit $p \colon C \to S$ une courbe stable de genre un et $H \subset C^{\reg}$ un sous-sch\'ema fini localement libre sur $S$ de l'ouvert de lissit\'e de $C$. On v\'erifie fibre par fibre que $H$ est automatiquement un diviseur de Cartier sur $C$. On suppose que les fibres de $H$ sont non vides.

\begin{proposition}[1.2]
Il existe un et un seul $S$-sch\'ema $C^{\natural}$, muni de $u \colon C \to C^{\natural}$, tel que
\begin{enumerate}[label=(\alph*)]
\item $C^{\natural}$ est une courbe stable de genre un sur $S$;
\item pour tout point g\'eom\'etrique $s$ de $S$, $C^{\natural}_{\bar{s}}$ se d\'eduit de $C_{\bar{s}}$ en contractant en un point chacune des composantes irr\'eductibles de $C_{\bar{s}}$ qui ne rencontrent pas $H_{\bar{s}}$.
\end{enumerate}
\end{proposition}

\noindent\textbf{A. Existence.} Un argument de passage \`a la limite nous ram\`ene \`a supposer $S$ noeth\'erien. Pour tout $s$ et tout $n > 0$, on a $\CO(nH_s)$ ample. Par dualit\'e (I.2.2 et II.1.2), on a donc $H^1(C_s, \CO(nH_s)) = 0$. D'apr\`es EGA III 7.8, les $p_* \CO(nH)$ ($n > 0$) sont localement libres de formation compatible \`a tout changement de base. Le m\^eme \'enonc\'e vaut pour $n=0$ (II.1.6). Soit $\CG$ l'alg\`ebre homog\`ene
\[
\CG = \bigoplus_{n \geq 0} p_* \CO(nH).
\]
Pour $C_{\bar{s}}$ irr\'eductible, $\CO(H_{\bar{s}})$ est ample, et $\CG_{\bar{s}}$ engendr\'e par $\bigoplus_{n \leq A} H^0(C_s, \CO(nH_s))$, avec $A$ ind\'ependant de $s$.

Supposons que $C_{\bar{s}}$ soit un polygone de N\'eron, et soit $\bar{C}_{\bar{s}}$ le polygone de N\'eron qui s'en d\'eduit en contractant en un point chaque composante irr\'eductible qui ne rencontre pas $H_{\bar{s}}$. On v\'erifie que $\CG_{\bar{s}}$ est l'alg\`ebre homog\`ene correspondante. Sur $\bar{C}_{\bar{s}}$, $\CO(n\bar{H}_{\bar{s}})$ est ample, et $\CG_{\bar{s}}$ est engendr\'e par $\bigoplus_{n \leq B} H^0(\bar{C}_s, \CO(n\bar{H}_s))$ avec $B$ ind\'ependant de $s$.

L'alg\`ebre homog\`ene $\CG = \bigoplus_{n \geq 0} p_* \CO_{C}(nH)$ est donc de type fini, et son spectre homog\`ene $C^{\natural}$ est propre et plat sur $S$, de formation compatible \`a tout changement de base. Le morphisme naturel $u \colon C \to C^{\natural}$ v\'erifie (a) et (b).

\medskip

\noindent\textbf{B. Unicit\'e.} Soit $u \colon C \to C^{\natural}$ v\'erifiant (i) et (ii). On v\'erifie fibre par fibre que $u$ induit un isomorphisme $u^{-1}(C^{\natural \reg}) \xrightarrow{\sim} C^{\reg}$ (EGA IV 17.8.2). Les arguments de A. montrent que $\CO(H)$ est ample et que $\CG = \bigoplus p_* \CO_{C}(nH)$ est plat et commute \`a tout changement de base. On v\'erifie fibre par fibre que $\CG \xrightarrow{\sim} \CG^{\natural}$, et $C^{\natural} = \Proj(\CG^{\natural})$.

\begin{proposition}[1.3]
Soit $p \colon C \to S$ une courbe elliptique g\'en\'eralis\'ee dont les fibres g\'eom\'etriques sont lisses ou sont des $nm$-gones. Il existe une et une seule courbe elliptique g\'en\'eralis\'ee $p \colon \bar{C} \to S$ munie de $u \colon C \to \bar{C}$, telle que
\begin{enumerate}[label=(\roman*)]
\item Les fibres g\'eom\'etriques de $\bar{C}$ sont lisses ou des $n$-gones. $\bar{C}_{\bar{s}}$ se d\'eduit de $C_{\bar{s}}$ en contractant en un point les composantes irr\'eductibles de $C_{\bar{s}}$ adh\'erence de celles des composantes connexes de $C^{\reg}_{\bar{s}}$ dont l'ordre dans $\pi_0(C^{\reg}_{\bar{s}})$ ne divise pas $n$.
\item Le diagramme suivant est commutatif:
\[
\xymatrix{
C^{\reg} \times_S C \ar[r] \ar[d] & C \ar[d]^{u} \\
\bar{C}^{\reg} \times_S \bar{C} \ar[r] & \bar{C}
}
\]
\end{enumerate}
\end{proposition}

Localement pour la topologie \fppf\ de $S$, il existe $H \subset C$ comme en 1.1, qui rencontre exactement celles des composantes irr\'eductibles des $C_{\bar{s}}$ qu'on ne d\'esire pas contracter (on peut prendre $H = C_n$, fini et plat d'apr\`es (II.1.20)). D'apr\`es 1.2, il existe une et une seule courbe stable de genre un $\bar{C}$ munie de $u \colon C \to \bar{C}$, qui v\'erifie (i). Le sous-groupe $u^{-1}(\bar{C}^{\reg})$ de $C^{\reg}$ agit sur $\bar{C}$ par transport de structure; puisque $u^{-1}(\bar{C}^{\reg}) \simeq C^{\reg}$, ceci fournit la structure de courbe elliptique g\'en\'eralis\'ee voulue sur $\bar{C}$.

\begin{example}[1.4]
Soit $C/S$ une courbe elliptique g\'en\'eralis\'ee, et appliquons 1.3 avec $n=1$. Localement sur $S$, l'hypoth\`ese de 1.3 est v\'erifi\'ee, pour $m$ convenable (II.1.15). Il existe donc une et une seule courbe elliptique g\'en\'eralis\'ee $c(C)$ sur $S$, \`a fibres g\'eom\'etriques irr\'eductibles, munie de $u \colon C \to c(C)$ et telle que $u$ induise un isomorphisme
\begin{equation}\label{eq:1.4.1}
c(C)^{\reg} \xrightarrow{\sim} C^{\reg}.
\end{equation}

D'apr\`es la d\'emonstration de 1.3, la courbe $c(C)$ est en effet uniquement d\'etermin\'ee par la condition (i) de 1.3, qui r\'esulte de (1.4.1), tandis que la structure de courbe g\'en\'eralis\'ee est uniquement d\'etermin\'ee par la section neutre (II.2.7), implicite dans 1.4.1.
\end{example}

\subsubsection{}
Soit $S$ le spectre d'un anneau de valuation discr\`ete complet \`a corps r\'esiduel alg\'ebriquement clos, et $\eta, s$ ses points g\'en\'eriques et sp\'eciaux. Pour toute extension finie $k(\eta')$ de $k(\eta)$, nous noterons $(S', \eta', s')$ le trait normalis\'e de $S$ dans $k(\eta')$. Soit $C_{\eta}$ une courbe elliptique sur $\eta$. Nous disons que $C_{\eta}$ \`a r\'eduction stable si le mod\`ele minimal de $C_{\eta}$ sur $S$ est une courbe stable de genre un.

\begin{proposition}[1.6]
\begin{enumerate}[label=(\roman*)]
\item Il existe une extension finie $k(\eta')$ de $k(\eta)$ telle que la courbe $C_{\eta'} = C_{\eta} \otimes_{\eta} \eta'$ ait r\'eduction stable sur $S'$.
\item Si $C_{\eta}$ \`a r\'eduction stable, le mod\`ele minimal $C$ de $C_{\eta}$ sur $S$ admet une unique structure de courbe elliptique g\'en\'eralis\'ee prolongeant celle de $C_{\eta}$.
\item Si de plus les points d'ordre $n$ de $C_{\eta}$ sont d\'efinis sur $\bar{\eta}$, soit le mod\'ele minimal $C$ de $C$ est lisse, soit $C_s$ est un $m$-gone, avec $n \mid m$, et il existe une courbe elliptique g\'en\'eralis\'ee $\bar{C}$ sur $S$, de fibre sp\'eciale lisse ou un $n$-gone et un isomorphisme $\alpha$ de sa fibre g\'en\'erique avec $C_{\eta}$.
\item Le couple $(\bar{C}, \alpha)$ est unique \`a isomorphisme unique pr\'es.
\end{enumerate}
\end{proposition}

\noindent\textit{Preuve.} (i). C'est un th\'eor\`eme de N\'eron et Kodaira ([4] ou [20] ou [14]).

(ii). La propri\'et\'e universelle du mod\`ele de N\'eron assure que $e_{\eta}$ se prolonge en une section de $C^{\reg}$ sur $S$. Si $u \in C^{\reg}(S)$, la translation $u+$ de $C_{\eta}$ se prolonge par transport de structure en un automorphisme de $C$; de m\^eme, $x \mapsto -x$ se prolonge \`a $C$. Soit $U$ le plus grand ouvert de $C$ tel que la loi $+$ de $C_{\eta}$ se prolonge en $+ \colon U \times_S C \to C$. Il suffit de prouver que $U = C^{\reg}$: ceci acquis, les identit\'es exprimant que $+$ est une structure de courbe elliptique g\'en\'eralis\'ee r\'esulteront de la densit\'e sch\'ematique de $C_{\eta}$ dans $C$. Pour $v \in C^{\reg}(S)$, $U$ est stable par translation par $v$: on prolonge $+$ \`a $v+U$ par $(v+u)+x = v+(u+x)$. Puisque $U$ est aussi stable par localisation \'etale sur $S$ (on le v\'erifie par un argument de descente; rappelons que la formation du mod\`ele minimal est compatible \`a la localisation \'etale sur $S$), il suffit de montrer que $U_s \neq \emptyset$.

Soit $x$ un point maximal (= g\'en\'erique) de $C_s$, et soit $S_1$ le trait spectre de l'anneau local de $C$ en $x$. Puisque $C$ est lisse sur $S$ en $x$, $C_1 = C \times_S S_1$ est encore un mod\`ele minimal sur $S_1$. Soit $u \in C_1(S_1)$ d\'efini par l'inclusion $\CO_{S_1} \subset \CO_{C_1}$. Par transport de structure, la translation $u+$ de la fibre g\'en\'erique de $C_1$ se prolonge en un automorphisme de $C_1$, d'o\`u
\[
+ \colon C_{1\eta} \times_{S_1} C_1 \longrightarrow C_1
\]
se prolonge en $+ \colon U_1 \times_{S_1} C_1 \to C_1$, avec $x \in U_1$.

(iii). On sait que $C_n$ est quasi-fini et plat. Puisque d'apr\`es l'hypoth\`ese et la propri\'et\'e universelle du mod\'ele de N\'eron, tous les points g\'eom\'etriques de $(C_{\bar{\eta}})_n$ proviennent de sections de $C_n$ sur $S$, ce groupe est m\^eme fini sur $S$ et $(C_s)_n$ est de rang $n^2$. La fibre sp\'eciale $C_s$ est donc lisse ou un $m$-gone, avec $n \mid m$ (cf.~II.1.13), et on applique 1.3.

(iv). Soit $\bar{C}$ une courbe elliptique g\'en\'eralis\'ee sur $S$, de fibre g\'en\'erale $C_{\eta}$ et telle que $\bar{C}_s$ soit lisse ou un polygone \`a $n$ c\^ot\'es. En un point de non lissit\'e $u$ de $\bar{C}$, $\bar{C}$ est formellement isomorphe au compl\'et\'e \`a l'origine de la courbe affine dans $\Af^2_S$ d'\'equation $xy = t^k$, pour $k$ convenable. Par homog\'en\'eit\'e, $k$ ne d\'epend pas du point choisi. Pour $k > 1$, $\bar{C}$ est singulier en $u$. Ces singularit\'es se r\'esolvent en \'eclatant de fa\c{c}on it\'er\'ee les points singuliers, ce processus aboutissant \`a remplacer $u$ par $k-1$ droites projectives:
\[
\xymatrix{
\tilde{C} \ar[d] \\ \bar{C}
}
\quad
\xymatrix{
A_k \ar[r] \ar[d] & \tilde{C} \ar[d] \\ A_1 \ar[r] & \bar{C}
}
\]
(cf.~[8]). Soit $\tilde{C}$ la courbe r\'esolue. Il n'existe pas de courbe exceptionnelle de premi\`ere esp\`ece contenue dans $\tilde{C}_s$, de sorte que $\tilde{C}$ est le mod\`ele minimal de $C_{\eta}$. La courbe $\bar{C}$ est lisse si et seulement si le mod\`ele minimal $\tilde{C}$ de $C_{\eta}$ est lisse, auquel cas $\bar{C} = \tilde{C}$. Sinon, $\bar{C}$ est un polygone de N\'eron \`a $nk$ c\^ot\'es, et $\tilde{C}$ s'en d\'eduit par contraction de la fa\c{c}on utilis\'ee dans la preuve de l'existence.
\qed


\subsection{Structures de niveau $n$}

\subsubsection{}
Soit $n$ un entier. Nous ne consid\'ererons dans ce \S\ que des sch\'emas $S$ sur lesquels $n$ est inversible.

Dans ce \S, nous noterons $\mll(n)$ le champ alg\'ebrique sur $\ZZ[1/n]$ qui classifie les courbes elliptiques g\'en\'eralis\'ees $C/S$, de fibres g\'eom\'etriques lisses ou des polygones de N\'eron \`a $n$ c\^ot\'es. D'apr\`es (II.1.15), $\mll(n)$ est un ouvert de $\mll^*$, donc un champ alg\'ebrique (III.2.5) lisse sur $\ZZ[1/n]$. Soit $f_n$ la section de $\mll(n)$ d\'efinie par la courbe II.1.9 sur $\ZZ[1/n]$.

\begin{proposition}[2.2]
Le champ alg\'ebrique $\mll(n)$ est propre et lisse sur $\ZZ[1/n]$. Son lieu \`a l'infini $\mll^{\infty}(n) = \mll(n) \cap \mll^{*,\infty}$ est l'image de $f_n$.
\end{proposition}

La seconde assertion n'est mise que pour rappel. Elle r\'esulte de III.2.5 et implique que $\mll^{\circ}(n)$ est dense dans $\mll(n)$. La propret\'e r\'esulte alors de 1.6 et du crit\`ere valuatif de propret\'e.

\subsubsection{}
Pour $C/S$ du type consid\'er\'e, $C_n$ est fini \'etale de rang $n^2$ (II.1.20), donc localement isomorphe \`a $(\ZZ/n)^2$ (pour la topologie \'etale). Une structure de niveau $n$ sur $C$ est un isomorphisme
\[
\alpha \colon C_n \xrightarrow{\sim} (\ZZ/n)^2.
\]

\begin{definition}[2.4]
$\mll_n[1/n]$ est le champ alg\'ebrique sur $\ZZ[1/n]$ suivant: pour $S$ un sch\'ema sur lequel $n$ est inversible, $\mll_n[1/n](S)$ est la cat\'egorie des courbes elliptiques g\'en\'eralis\'ees $C/S$, de fibres g\'eom\'etriques lisses ou des polygones de N\'eron \`a $n$ c\^ot\'es, munies d'une structure de niveau $n$.
\end{definition}

Il est clair que $\mll_n[1/n]$ est fini \'etale et repr\'esentable sur $\mll(n)$, et m\^eme un espace principal homog\`ene de groupe $\Gl(2, \ZZ/n)$. On note $\mll^{\infty}_n[1/n]$ l'image r\'eciproque de $\mll^{\infty}(n)$, image du sous-sch\'ema de non lissit\'e de la courbe universelle sur $\mll_n[1/n]$. Le th\'eor\`eme suivant r\'esulte aussit\^ot de 2.2.

\begin{theorem}[2.5]
$\CM_n[1/n]$ est propre et lisse sur $\ZZ[1/n]$ et $\mll^{\infty}_n[1/n]$ est fini et \'etale sur $\ZZ[1/n]$.
\end{theorem}

\subsubsection{}
Rappelons qu'un champ alg\'ebrique $\CG$ ``est'' un espace alg\'ebrique si les objets qu'il classifie n'ont pas d'automorphismes. Plus pr\'ecis\'ement, si pour tout corps alg\'ebriquement clos $k$, les objets de $\CG(k)$ n'ont pas d'automorphisme non trivial, alors
\begin{enumerate}[label=\alph*)]
\item pour tout sch\'ema $S$, les objets de $\CG(S)$ n'ont pas d'automorphisme non trivial;
\item le foncteur $S \mapsto ($l'ensemble des classes d'isomorphie dans $\CG(S))$ est repr\'esentable par un espace alg\'ebrique.
\end{enumerate}

\begin{theorem}[2.7]
Si $n \geq 3$, $\mll_n[1/n]$ est un espace alg\'ebrique.
\end{theorem}

Soient $C$ une courbe elliptique g\'en\'eralis\'ee sur un corps alg\'ebriquement clos $k$, lisse ou un polygone de N\'eron \`a $n$ c\^ot\'es. Il faut prouver qu'un automorphisme $\sigma$ de $C$, trivial sur $C_n$, est l'identit\'e. Pour $C$ une courbe elliptique, on applique ([22], app.\ \`a Exp.~17). Pour $C$ un polygone de N\'eron, on v\'erifie sur II.1.10 que $\Aut(C)$ agit fid\`element sur $C_n$.

\begin{variante}[2.8]
Cet argument, et III.2.1, impliquent aussit\^ot que le foncteur
\[
F_n \colon S \mapsto (\text{l'ensemble des classes d'isomorphie de courbes elliptiques g\'en\'eralis\'ees sur } S \text{ munies d'une structure de niveau } n)
\]
est un faisceau pour la topologie \fpqc. Les arguments de III.2.5 montrent alors que ce foncteur v\'erifie le crit\`ere de repr\'esentabilit\'e de M.~Artin. Pour prouver que $F_n$ est repr\'esentable par un espace alg\'ebrique ($n \geq 3$), il n'est donc pas n\'ecessaire de g\'en\'eraliser au pr\'alable le crit\`ere de M.~Artin au cas des champs alg\'ebriques.
\end{variante}

\begin{corollary}[2.9]
Si $n \geq 3$, $\CM_n[1/n]$ est un sch\'ema projectif et lisse sur $\ZZ[1/n]$.
\end{corollary}

Un espace alg\'ebrique r\'egulier en courbes sur $\ZZ$ est en effet toujours quasi-projectif, donc un sch\'ema. Voici une d\'emonstration plus constructive: par voie transcendante (cf.~\S 5), on v\'erifie qu'apr\`es extension des scalaires \`a $\CC$, $\mll^{\infty}_n$ rencontre chaque composante irr\'eductible de $\CM_n$. Par sp\'ecialisation, on en d\'eduit que le diviseur $\mll^{\infty}_n$ rencontre chaque composante irr\'eductible de chaque fibre g\'eom\'etrique de $\CM_n$ (utiliser que $\CM_n$ est propre et lisse). D\`es lors, $\CO(\mll^{\infty}_n)$ est un faisceau inversible relativement ample.


\subsection{Structures de niveau $H$}

\subsubsection{}
Soient $n$ un entier et $H$ un sous-groupe de $\Gl(2, \ZZ/n)$. Soient $S$ un sch\'ema sur lequel $n$ est inversible et $C$ une courbe elliptique sur $S$. Pour tout $S$-sch\'ema $T$, soit $F(T)$ l'ensemble des structures de niveau $n$ sur $C_T$, l'image r\'eciproque de $C$ sur $T$. Le foncteur $F$ est repr\'esent\'e par un rev\^etement fini \'etale de $S$ qui est un espace principal homog\`ene de groupe $\Gl(2, \ZZ/n)$.

Le faisceau pour la topologie \'etale engendr\'e par le pr\'efaisceau
\[
T \mapsto H \backslash F(T)
\]
($H$ agissant \`a gauche, il vaudrait peut-\^etre mieux \'ecrire $H \backslash F(T)$). Dans le langage canonique, ce foncteur est repr\'esent\'e par un rev\^etement fini \'etale de $S$.

Une structure de niveau $H$ sur $C$ est un \'el\'ement de $F_H(S)$. Un isomorphisme $\alpha \colon C_n \xrightarrow{\sim} (\ZZ/n)^2$ d\'efinit une structure de niveau $H$, et deux tels isomorphismes d\'efinissent la m\^eme structure si et seulement si, localement sur $S$, il existe $g \in H$ tel que $g\alpha = \beta$. Ce $g$ est uniquement d\'etermin\'e, donc constant sur chaque composante connexe de $S$. On prendra garde qu'une structure de niveau $H$ ne provient en g\'en\'eral pas d'un $\alpha \colon C_n \xrightarrow{\sim} (\ZZ/n)^2$. Elle ne provient de tels $\alpha$ que localement pour la topologie \'etale. Par exemple, pour $H = \Gl(2, \ZZ/n)$, $C$ a une et une seule structure de niveau $H$; elle ne vient d'un $\alpha$ que si $C_n$ est isomorphe \`a $(\ZZ/n)^2$.

\begin{definition}[3.2]
$\mll_H[1/n]$ est le champ alg\'ebrique suivant: pour $S$ un sch\'ema sur lequel $n$ est inversible, $\mll_H[1/n](S)$ est la cat\'egorie des courbes elliptiques $C/S$, munies d'une structure de niveau $H$ (les morphismes sont les isomorphismes).
\end{definition}

Il est clair que $\mll_H[1/n]$ est fini \'etale et localement repr\'esentable sur $\mll_1[1/n]$.

\begin{definition}[3.3]
Le champ alg\'ebrique $\mll'_H$ est le normalis\'e de $\mll_1$ dans $\mll_H[1/n]$.
\end{definition}

\begin{theorem}[3.4] \ \begin{enumerate}[label=(\roman*)]
\item $\mll_H$ est propre sur $\Spec(\ZZ)$.
\item $\mll^h_H[1/n]$ est lisse sur $\Spec(\ZZ[1/n])$, et $\mll^{\infty}_H[1/n]$ est le compl\'ement d'un sous-champ $\mll^{h,\infty}_H[1/n]$ fini et \'etale sur $\Spec(\ZZ[1/n])$.
\end{enumerate}
\end{theorem}

\noindent\textit{Preuve.} (i) Par d\'efinition, $\mll_H$ est fini et relativement repr\'esentable sur $\mll_1$. Sa propret\'e r\'esulte donc de celle de $\mll_1$.

(ii) L'assertion r\'esulte du lemme d'Abhyankar (SGA 1, XIII 5) applicable ici parce que $\mll_1[1/n]$ est lisse sur $\ZZ[1/n]$, que $\mll^{\infty}_1[1/n]$ est un diviseur de $\mll_1[1/n]$, lisse sur $\ZZ[1/n]$, et dont le point g\'en\'erique est de caract\'eristique $0$, et que $\mll_H[1/n]$ est fini \'etale sur $\mll_1[1/n] - \mll^{\infty}_1[1/n]$.

En 6.7, nous donnerons de 3.4 une d\'emonstration ``modulaire''.

Pour $H = \{e\}$, nous poserons $\mll_n = \mll_H$. Cette notation est l\'egitim\'ee par la proposition suivante.

\begin{proposition}[3.5]
Pour $H = \{e\}$, $\mll_H[1/n]$ est le champ $\mll_n[1/n]$ de 2.4.
\end{proposition}

La construction 1.4 d\'efinit un morphisme de champs
\[
c \colon \mll_n[1/n] \longrightarrow \mll_1[1/n] \colon (C, \alpha) \mapsto c(C).
\]
Puisque $\mll_1[1/n]$ est normal (m\^eme lisse sur $\ZZ[1/n]$), propre sur $\ZZ[1/n]$, et que clairement $\mll^{\circ}_n[1/n] \xrightarrow{\sim} \mll_H[1/n]$, il nous suffit de prouver que $c$ est fini et localement repr\'esentable. Il est clair que les fibres de $c$ sont finies. Pour tester la repr\'esentabilit\'e locale, il suffit d\`es lors de v\'erifier (cf.~2.6):

\begin{lemma}[3.5.1]
Pour $k$ un corps alg\'ebriquement clos, $C$ un objet de $\mll_n[1/n](k)$ et $c(C)$ son image dans $\mll_1(k)$, l'application naturelle $i \colon \Aut(C) \to \Aut(c(C))$ est injective.
\end{lemma}

Pour $n \geq 3$, $\Aut(C)$ est trivial (cf.~2.7). Pour $n = 2$, $\Aut(C) = \{\pm 1\}$. Enfin, pour $n=1$, $i$ est l'identit\'e.

\subsubsection{}
Soient $n$ et $m$ deux entiers, avec $n \mid m$. Soient $H \subset \Gl(2, \ZZ/n)$, et $H'$ son image r\'eciproque dans $\Gl(2, \ZZ/nm)$. Pour $E$ une courbe elliptique sur $S$, une structure de niveau $E_{nm} \xrightarrow{\sim} (\ZZ/nm)^2$ d\'efinit une structure de niveau $n$, rendant commutatif le diagramme
\[
\xymatrix{
E_{nm} \ar[r]^{\alpha'} \ar[d]_{\times m} & (\ZZ/nm)^2 \ar[d] \\
E_n \ar[r]^{\alpha} & (\ZZ/n)^2
}
\]
Pour $nm$ inversible sur $S$, cette construction d\'efinit un isomorphisme du faisceau des $H'$-structures sur $E$ avec celui des $H$-structures. On a
\[
\mll_{H'}[1/nm] \simeq \mll_H[1/n] \otimes_{\ZZ} \ZZ[1/m].
\]
Puisque $\mll_H[1/n]$ est fini et \'etale sur $\mll^h_1[1/n]$, c'est le normalis\'e de $\mll^h_1[1/n]$ dans $\mll_{H'}[1/nm]$. On en d\'eduit un isomorphisme
\[
\mll_{H'} \simeq \mll_H.
\]
Le champ $\mll_H$ ne d\'epend donc que du sous-groupe ouvert $K$ de $\Gl(2, \hat{\ZZ})$ image r\'eciproque de $H$. Nous le noterons parfois $\mll_K$.

\subsubsection{}
Soient $n$ et $m$ deux entiers premiers entre eux, $H \subset \Gl(2, \ZZ/n)$ et $I \subset \Gl(2, \ZZ/m)$. Soient $K$ et $L$ les images r\'eciproques de $H$ et $I$ dans $\Gl(2, \hat{\ZZ})$. Les morphismes de r\'eduction mod $n$ et $m$ induisent un isomorphisme $\Gl(2, \ZZ/nm) \simeq \Gl(2, \ZZ/n) \times \Gl(2, \ZZ/m)$. Via cet isomorphisme, $K \cap L$ est l'image r\'eciproque de $H \times I \subset \Gl(2, \ZZ/n) \times \Gl(2, \ZZ/m)$.

\begin{construction}[3.8]
Avec les notations pr\'ec\'edentes,
\[
\mll_{K \cap L}[1/nm] \simeq \mll_K[1/n] \times_{\mll_1} \mll_L[1/m].
\]
\end{construction}

Soit $C/S$ une courbe elliptique. Les applications $x \mapsto x^m$ et $x \mapsto x^n$ d\'efinissent un isomorphisme
\begin{equation}\label{eq:3.8.1}
C_{nm} \xrightarrow{\sim} C_n \times C_m.
\end{equation}
Si $\alpha$ (resp.\ $\beta$) est une structure de niveau $H$ (resp.\ $I$) sur $C$, localement repr\'esent\'ee par $\alpha \colon C_n \xrightarrow{\sim} (\ZZ/n)^2$, on d\'efinit la structure $\alpha \times \beta$ de niveau $H \times I$ sur $C$ comme \'etant la classe de l'isomorphisme $\alpha \times \beta$ rendant commutatif le diagramme
\[
\xymatrix{
C_{nm} \ar[r]^{\alpha \times \beta} \ar[d]_{\eqref{eq:3.8.1}} & (\ZZ/nm)^2 \ar[d]^{\text{r\'eduction}} \\
C_n \times C_m \ar[r]^{\alpha \times \beta} & (\ZZ/n)^2 \times (\ZZ/m)^2
}
\]

L'isomorphisme 3.8 est d\'efini par $(C, \alpha, \beta) \mapsto (C, \alpha \times \beta)$.

\begin{proposition}[3.9]
Avec les notations de 3.7, l'isomorphisme 3.8 se prolonge en
\begin{equation}\label{eq:3.9.1}
\mll_{K \cap L} \simeq \mll_K \times_{\mll_1} \mll_L.
\end{equation}
\end{proposition}

\noindent\textit{Preuve.} Il suffit de prouver que $\mll_K \times_{\mll_1} \mll_L$ est normal. Prouvons-le au-dessus de $\ZZ[1/m]$. Au-dessus de $\ZZ[1/n]$, la d\'emonstration est sym\'etrique. $\mll_K$ est lisse sur $\ZZ[1/n]$, donc plat sur $\mll_1$. $\mll_L[1/m]$ est donc plat sur $\mll_1 \otimes \ZZ[1/m]$ et de Cohen-Macaulay (I.7.2 et 7.1). $\mll_L[1/m]$ est m\^eme \'etale sur $\mll_1$, donc $\mll_K \times_{\mll_1} \mll_L[1/m]$, \'etale sur $\mll_K$, est normal.

D'apr\`es le lemme d'Abhyankar, la ramification de $\mll_H[1/n]$ sur $\mll_1[1/n]$ (resp.\ $\mll_I[1/m]$ sur $\mll_1[1/m]$) est mod\'er\'ee. Les indices de ramifications pour $\mll_K$ et $\mll_L$ sont donc premiers entre eux (en fait, des diviseurs de $n$ et $m$). Il en r\'esulte que $\mll_H[1/mn] \times_{\mll_1} \mll_I[1/mn]$ est lisse, et il ne reste plus qu'\`a appliquer le crit\`ere de Serre (I.7.2).
\qed

\begin{corollary}[3.9.2]
Si $n$ est le produit de deux entiers premiers entre eux et $\geq 3$, alors $\CM_n$ est un sch\'ema.
\end{corollary}

Soit $n = n' \cdot n''$. Le champ $\mll_n[1/n']$ est un sch\'ema, car repr\'esentable sur $\mll_{n'}$. De m\^eme, $\mll_n[1/n'']$, donc $\mll_n$, est un sch\'ema.

\begin{proposition}[3.10] \ \begin{enumerate}[label=(\roman*)]
\item L'espace grossier $\CM_H$ d\'efini par $\mll_H$ est propre et plat sur $\Spec(\ZZ)$.
\item C'est le normalis\'e dans $\mll_H[1/n]$ de $\CM_1$ (isomorphe \`a $\PP^1$, voir VI.1.1).
\item $\mll^{\circ}_H[1/n] \simeq \mll^{\circ}_1[1/n]$ est fini et \'etale sur $\Spec(\ZZ[1/n])$.
\item Les fibres g\'eom\'etriques de $\mll^h_H[1/n] \to \Spec(\ZZ[1/n])$ sont g\'eom\'etriquement unibranches.
\end{enumerate}
\end{proposition}

\noindent\textit{Preuve.} La propret\'e de $\CM_H$ r\'esulte de celle de $\mll_H$. De la description de l'henselis\'e strict d'un anneau local de $\CM_H$ donn\'ee en I.8.2 et de la normalit\'e de $\CM_H$, il r\'esulte que $\CM_H$ est normal. De m\^eme, $\CM_H$ est plat sur $\ZZ$. Puisque $\mll_H \to \mll_1$ donc $\CM_H \to \CM_1$ est fini, (ii) r\'esulte de (i) et de la normalit\'e.

Il suffit de prouver (iii) et (iv) lorsque $n \geq 3$ (cf.~3.6), auquel cas $\mll^h_H[1/n] = \mll^h_0[1/n]$. Avec les notations de [8] \S 4 on a alors
\[
\mll^{\circ}_n[1/n] \simeq \mll^{\circ}_1[1/n] \times_{\mll_1[1/n]} \mll_n[1/n],
\]
de sorte que $\mll^{\circ}_H[1/n] = \mll^{\circ}_1[1/n] \times_{\mll_1[1/n]} \mll_H[1/n] = \mll^{\circ}_1[1/n] / H$. Ceci \'etant, (iii) r\'esulte de (ii).

L'application
\[
(\CM_n \otimes \FF_p)/H \to (\CM_n/H) \otimes \FF_p = \CM_H \otimes \FF_p
\]
est radicielle. Si $p \nmid n$, $\CM_n \otimes \FF_p$ est lisse, donc $(\CM_n \otimes \FF_p)/H$ est lisse, et (iv) en r\'esulte.
\qed


\subsubsection{}
La cat\'egorie des courbes elliptiques \`a isog\'enie pr\`es sur un sch\'ema $S$ est la cat\'egorie d\'eduite de celle des courbes elliptiques en inversant les isog\'enies. Notant $E \sim E'$ la courbe elliptique \`a isog\'enie pr\`es sur $S$ sous-jacente \`a $E/S$, on a
\[
\Hom(E_0, E'_0) = \Hom(E,E') \otimes_{\ZZ} \QQ.
\]

\subsubsection{}
Pla\c{c}ons-nous sur un corps alg\'ebriquement clos de caract\'eristique $0$, $k$. Pour toute courbe elliptique $E$ sur $k$, soit
\[
T(E) = \varprojlim_n E_n \qquad (\text{isomorphe } \`a \hat{\ZZ}^2).
\]
Le groupe des points de division de $E$ est canoniquement isomorphe \`a $V(E)/T(E)$ (\`a $x \in E_n$ on associe l'image dans $E_n$ de $nx \in T(E)$).

Une isog\'enie $E \to F$ induit un isomorphisme $V(E) \xrightarrow{\sim} V(F)$. Ceci permet de d\'efinir $V(E_0)$ pour $E_0$ une courbe elliptique \`a isog\'enie pr\`es. Si $f \colon E \to F$ est une isog\'enie, on a
\[
\Ker(f) = T(F)/T(E) \subset V(E)/T(E).
\]
L'application $f \mapsto T(F) \subset V(E)$ identifie les isog\'enies de source $E$ aux ``r\'eseaux'' $T$ contenant $T(E)$. On en d\'eduit que le foncteur
\[
(\text{courbes elliptiques}) \longrightarrow (\text{courbes elliptiques \`a isog\'enie pr\`es } E_0, \text{ munies d'un ``r\'eseau'' } T \subset V(E_0))
\]
est une \'equivalence de cat\'egories.

\subsubsection{}
Soit $E_0$ une courbe elliptique \`a isog\'enie pr\`es sur $k$. Un r\'eseau $T \subset V(E_0)$ peut s'interpr\'eter comme une classe $\bmod\ \Gl(2, \hat{\ZZ})$ d'isomorphismes
\[
\delta \colon V(E_0) \xrightarrow{\sim} \Af_f^2.
\]
\`A $\delta$ on associe $\delta^{-1}(\hat{\ZZ}^2)$.

Soit $K \subset \Gl(2, \Af_f)$ l'image r\'eciproque de $H \subset \Gl(2, \ZZ/n)$. A une courbe $E$, \`a une structure de niveau $H$ sur $E$ on associe l'ensemble des $\delta \colon V(E) \xrightarrow{\sim} \Af_f^2$ tels que $\delta^{-1}(\hat{\ZZ}^2) = T$ et que
\[
\beta_n \colon E_n = T/nT \xrightarrow{\sim} \hat{\ZZ}^2/n\hat{\ZZ}^2 = (\ZZ/n)^2
\]
d\'efinisse la structure de niveau $H$. Les $\beta$ forment une classe lat\'erale sous $K$, et l'ensemble des courbes elliptiques munies d'une structure de niveau $H$, $(E, \alpha)$, avec $E \sim E' = E_0$, s'identifie \`a
\[
K \backslash \Isom(V(E_0), \Af_f^2).
\]

\subsubsection{}
Soit $g \in \Gl(2, \Af_f)$ tel que $gKg^{-1} \subset \Gl(2, \hat{\ZZ})$. Si on identifie une courbe elliptique avec structure de niveau $K$ (= de niveau $H$) \`a $(E_0, \delta)$, avec $\delta \in K \backslash \Isom(V(E_0), \Af_f^2)$, l'application $(E_0, \delta) \mapsto (E_0, g\delta)$ associe \`a une courbe de niveau $K$ une courbe de niveau $gKg^{-1}$.

Ces constructions gardent un sens sur un sch\'ema de caract\'eristique $0$ quelconque et d\'efinissent
\begin{equation}\label{eq:3.14.1}
c_g \colon \mll_K \otimes \QQ \longrightarrow \mll_{gKg^{-1}} \otimes \QQ
\end{equation}
et un isomorphisme $[g]$ entre l'image r\'eciproque par $c_g$ de la courbe universelle sur $\mll_{gKg^{-1}} \otimes \QQ$ et la courbe universelle sur $\mll_K \otimes \QQ$, \`a isog\'enie pr\`es.

Par passage \`a la limite sur $K$, on en d\'eduit une action de $\Gl(2, \Af_f)$ sur le sch\'ema $\varprojlim_K \mll_K \otimes \QQ$, et sur la courbe elliptique universelle \`a isog\'enie pr\`es.

\subsubsection{}
Explicitons l'application compos\'ee
\[
\mll_K \otimes \QQ \xrightarrow{c_g} \mll_{gKg^{-1}} \otimes \QQ \xrightarrow{j} \mll_1 \otimes \QQ = \PP^1_{\QQ}.
\]
Prenons $n$, et soit $m$ tels que
\[
g(\hat{\ZZ}^2) \subset \frac{1}{n}\hat{\ZZ}^2 \quad \text{et} \quad g^{-1}(\hat{\ZZ}^2) \cap \hat{\ZZ}^2 \supset m\hat{\ZZ}^2.
\]
Soit $B$ le sous-groupe de $(\frac{1}{n}\ZZ/\ZZ)^2$ image de $\hat{\ZZ}^2$ par la multiplication par $n$, $\hat{\ZZ}^2 \to \frac{1}{n}\hat{\ZZ}^2/\hat{\ZZ}^2 \simeq (\ZZ/n)^2$.

\`A une courbe elliptique avec structure de niveau $H$, $(E, \alpha)$, on associe la courbe elliptique $E/\alpha^{-1}(B)$, pour $\alpha \colon E_n \xrightarrow{\sim} (\ZZ/n)^2$ un quelconque repr\'esentant de $\alpha$. L'isog\'enie $[g]$ est $n/m$ fois la projection $E \to E/\alpha^{-1}(B)$.

\begin{proposition}[3.16]
L'isomorphisme 3.14.1 se prolonge en un isomorphisme
\[
c_g \colon \mll_K \xrightarrow{\sim} \mll_{gKg^{-1}}.
\]
L'isomorphisme $[g]$ se prolonge en un isomorphisme
\[
[g] \colon c_g^*(\text{courbe elliptique sur } \mll_{gKg^{-1}}) \xrightarrow{\sim} (\text{courbe elliptique sur } \mll_K) \otimes \QQ.
\]
\end{proposition}

\noindent\textit{Preuve.} Montrons tout d'abord que (3.15.1) se prolonge en $c_g \colon \mll_K \to \mll_{gKg^{-1}}$.

Soit $b$ l'ordre de $B$ comme en 3.15 et soit $\CI$ le champ relativement repr\'esentable sur $\mll_1$ qui repr\'esente le foncteur des sous-sch\'emas en groupes localement libres de rang $b$ de la courbe elliptique universelle $E$. La th\'eorie des sch\'emas de Hilbert montre que la projection $\CI \to \mll_1$ est un morphisme propre et repr\'esentable. En vertu du lemme suivant, il est fini.

\begin{lemma}[3.17]
Soient $E$ une courbe elliptique sur $k$ alg\'ebriquement clos et $b$ un entier. Il n'y a qu'un nombre fini de sous-sch\'emas en groupes de rang $b$ de $E$.
\end{lemma}

\noindent\textit{Preuve de 3.17.} On proc\'ede par r\'ecurrence sur $b$.
\begin{enumerate}[label=\alph*)]
\item Il n'y a qu'un nombre fini de sous-sch\'emas \'etales de rang $b$ de $E$. Ils s'identifient aux sous-groupes d'ordre $b$ de $E_b(k)$.
\item Si $H \subset E$ n'est pas \'etale, $k$ est de caract\'eristique $p > 0$ et $H \supset \Ker(F)$. Les sous-groupes non \'etales de rang $b$ de $E$ correspondent donc aux sous-groupes de rang $b/p$ de $E^{(p)} = E/\Ker(F)$, et on conclut par l'hypoth\`ese de r\'ecurrence.
\qed
\end{enumerate}

\noindent\textit{Preuve de 3.16 (suite).} La construction 3.14 d\'efinit une section $s$ sur $\mll_K \otimes \QQ$ de $\CI \times_{\mll_1} \mll_K$.

Il r\'esulte du Main Theorem de Zariski et de la normalit\'e de $\mll_K$ que cette section se prolonge sur $\mll_K$, d\'efinissant un sous-groupe $\bar{B}$ de l'image inverse $E_0$ par $c$ de la courbe elliptique universelle sur $\mll_1$; le prolongement $c_g$ voulu est d\'efini par $E/\bar{B}$; l'isomorphisme $[g]$ est d\'efini comme en 3.15.

Le morphisme $c_g$ ainsi d\'efini est quasi-fini et repr\'esentable. Il identifie donc $\mll_K$ \`a un ouvert du normalis\'e de $\mll_{gKg^{-1}}$ dans $\mll_K \otimes \QQ \simeq \mll_{gKg^{-1}} \otimes \QQ$, d'o\`u
\[
c_g \colon \mll_K \xrightarrow{\sim} \mll_{gKg^{-1}}.
\]
Ce morphisme est un isomorphisme, d'inverse $c_{g^{-1}}$.

\subsubsection{}
On peut montrer que $c_g$ se prolonge en un isomorphisme
\[
\mll_K \xrightarrow{\sim} \mll_{gKg^{-1}}.
\]
Nous nous contenterons de montrer le r\'esultat analogue pour les sch\'emas grossiers.

\begin{proposition}[3.19]
Le morphisme de sch\'emas grossiers d\'eduit de 3.16 se prolonge en
\[
c_g \colon \CM_K \xrightarrow{\sim} \CM_{gKg^{-1}}.
\]
\end{proposition}

Soit $\Gamma \subset \CM_K \times \CM_{gKg^{-1}}$ l'adh\'erence du graphe $G$ de $c_g$. On a, ensemblistement,
\[
\Gamma = G \cup (\CM^{\infty}_K \times \CM^{\infty}_{gKg^{-1}}).
\]
Des lors, $\Gamma$ est fini sur $\CM_K$ et $\CM_{gKg^{-1}}$ et, d'apr\`es le Main theorem et 3.10, $\Gamma \xrightarrow{\sim} \CM_K$ et $\Gamma \xrightarrow{\sim} \CM_{gKg^{-1}}$ sont des isomorphismes. 3.19 en r\'esulte.

\subsubsection{}
Soit $C$ une courbe elliptique sur $S$, avec $n$ inversible sur $S$. Il est bien connu que le ``$e_n$-pairing'' d\'efinit un isomorphisme
\begin{equation}\label{eq:3.20.0}
e_n \colon C_n \otimes C_n \xrightarrow{\sim} \mu_n.
\end{equation}
Une structure de niveau $n$, $\alpha \colon C_n \xrightarrow{\sim} (\ZZ/n)^2$, d\'efinit donc un isomorphisme $\det(\alpha)^{-1} \colon \ZZ/n \xrightarrow{\sim} \mu_n$, i.e.\ une racine primitive $n$-i\`eme de l'unit\'e $\det(\alpha)^{-1}(1)$ sur $S$.

Notons $\ZZ[\zeta_n]$ l'anneau des entiers du corps des racines $n$-i\`emes de l'unit\'e. La construction $\alpha \mapsto \det(\alpha)$ d\'efinit un morphisme
\begin{equation}\label{eq:3.20.1}
d \colon \mll_n[1/n] \longrightarrow \Spec(\ZZ[\zeta_n]).
\end{equation}
Ce morphisme se prolonge au normalis\'e et d\'efinit
\begin{equation}\label{eq:3.20.2}
d \colon \CM_n \longrightarrow \Spec(\ZZ[\zeta_n]).
\end{equation}

De m\^eme, une structure de niveau $H$, $\alpha$, d\'efinit
\begin{equation}\label{eq:3.20.3}
d \colon \mll_H[1/n] \longrightarrow \Spec(\ZZ[\zeta_n]^{\det H}).
\end{equation}

Le groupe $(\ZZ/n)^*$ agit sur $\ZZ[\zeta_n]$; si $\ZZ[\zeta_n]^{\det H}$ est le sous-anneau des invariants de $\ZZ[\zeta_n]$ sous l'action de $\det H \subset (\ZZ/n)^*$, on trouve comme plus haut:
\begin{equation}\label{eq:3.20.4}
d \colon \CM_H \longrightarrow \Spec(\ZZ[\zeta_n]^{\det H})
\end{equation}
d\'efinissant
\[
d \colon \CM_H \longrightarrow \Spec(\ZZ[\zeta_n]^{\det H}).
\]

\subsubsection{}
Pla\c{c}ons-nous sur $\mll(m)[1/m]$ avec $n \mid m$. Si $C$ est la courbe elliptique g\'en\'eralis\'ee universelle, le $e_n$-pairing se prolonge en un isomorphisme
\[
e_n \colon C_n \otimes C_n \xrightarrow{\sim} \mu_n.
\]
En particulier, pour le polygone de N\'eron \`a $m$ c\^ot\'es standard, on trouve un isomorphisme
\begin{equation}\label{eq:3.21.1}
e_n \colon C_n \otimes C_n \xrightarrow{\sim} \mu_n.
\end{equation}
Identifions $\ZZ/n$ \`a $(\ZZ/m)_n$ par $i \mapsto (m/n)i$. On a
\begin{equation}\label{eq:3.21.2}
C_n \simeq C_m \times_{(\ZZ/m)^2} (\ZZ/n)^2.
\end{equation}
Selon les conventions de signe utilis\'ees pour d\'efinir $e_n$, (3.21.1) et (3.21.2) sont \'egaux ou oppos\'es.


\subsection{Exemples}

\subsubsection{}
Nous noterons $\Gamma_0(n)$, $\Gamma_{00}(n)$, $\Gamma'_{00}(n)$ et $\Gamma(n)$ les sous-groupes de $\Gl(2, \ZZ)$ form\'es des matrices $\begin{pmatrix} a & b \\ c & d \end{pmatrix}$ telles que
\begin{align*}
\Gamma_0(n) & \quad c \equiv 0 \pmod n; \\
\Gamma_{00}(n) & \quad c \equiv 0 \pmod n, \quad a \equiv 1 \pmod n; \\
\Gamma'_{00}(n) & \quad c \equiv 0 \pmod n, \quad a \equiv d \equiv 1 \pmod n; \\
\Gamma(n) & \quad \begin{pmatrix} a & b \\ c & d \end{pmatrix} \equiv \begin{pmatrix} 1 & 0 \\ 0 & 1 \end{pmatrix} \pmod n.
\end{align*}
Ces sous-groupes sont chacun image r\'eciproque de sous-groupes $H$ de $\Gl(2, \ZZ/n)$.

\subsubsection{}
Pour $\Gamma_0(n)$, $H$ est le sous-groupe de $\Gl(2, \ZZ/n)$ qui stabilise le sous-groupe $\ZZ/n \times \{0\}$ de $(\ZZ/n)^2$. Si $C$ est une courbe elliptique sur $S$, et que $\alpha$ est une structure de niveau $H$ sur $C$, repr\'esent\'ee par $\alpha \colon C_n \xrightarrow{\sim} (\ZZ/n)^2$, le sous-sch\'ema en groupes $\alpha^{-1}(\ZZ/n \times \{0\})$ de $C$ ne d\'epend donc que de $\alpha$. Si $n$ est inversible sur $S$, cette construction \'etablit un isomorphisme du faisceau des structures de niveau $H$ sur $C$ avec le faisceau des sous-sch\'emas en groupes de $C$ localement isomorphes \`a $\ZZ/n$. En d'autres termes:

\begin{construction}[4.3]
$\mll_0(n)[1/n]$ s'identifie au champ alg\'ebrique classifiant les courbes elliptiques $C/S$, avec $n$ inversible sur $S$, munies d'un sous-sch\'ema en groupes $A$ localement isomorphe \`a $\ZZ/n$.
\end{construction}

\subsubsection{}
Pour $C/S$ et $A$ comme en 4.3, le sous-groupe $C_n/A$ de $C/A$ est localement isomorphe \`a $\ZZ/n$, d'o\`u un morphisme
\[
w \colon (C,A) \longmapsto w(C,A) = (C/A, C_n/A).
\]
La multiplication par $n$ se factorise par un isomorphisme $C/C_n \xrightarrow{\sim} C$; celui-ci d\'efinit un isomorphisme
\[
ww(C,A) = ((C/A)/(C_n/A), (C/A)_n/(C_n/A)) \simeq (C/C_n, \ldots) \simeq (C,A)
\]
d'o\`u un isomorphisme $w^2 \simeq \Id$: $w$ est une involution. On v\'erifie que $w$ est du type 3.16, pour $g = \begin{pmatrix} 0 & -1 \\ n & 0 \end{pmatrix}$.

\subsubsection{}
Voici une description plus sym\'etrique du champ et de $w$. Associons \`a $(E,A)$ l'isog\'enie $E \to E/A$. On trouve que $\mll_0(n)[1/n]$ est encore le champ classifiant les isog\'enies
\[
E_1 \xrightarrow{\alpha} E_2
\]
de noyau cyclique d'ordre $n$.

Si $(E,A)$ d\'efinit $E_1 \xrightarrow{\alpha} E_2$, $w(E,A)$ d\'efinit une isog\'enie $E_2 \to E_1/(E_1)_n$, qui, via l'isomorphisme 4.4 de $E_1/(E_1)_n$ avec $E_1$, s'identifie au dual de Cartier de $\alpha$ dans le langage des isog\'enies.

\subsubsection{}
Soient $C/S$ et $A$ comme en 4.3. Pour chaque facteur premier $p$ de $n$, soit $A_p$ la composante $p$-primaire de $A$: $A = \prod_{p \mid n} A_p$, et soit $n(p)$ l'ordre de $A_p$, $n = \prod_{p \mid n} n(p)$. On pose $w_p(E,A) = (E/A_p, E_{pn(p)}/A_p)$. Les $w_p$ d\'efinissent des involutions de $\mll_0(n)[1/n]$, celles-ci commutent deux \`a deux, et $w$ est leur compos\'ee.

\subsubsection{}
Le sous-groupe $\Gamma_{00}(n)$ de $\Gl(2, \ZZ)$ est l'image r\'eciproque du sous-groupe $H$ de $\Gl(2, \ZZ/n)$ qui fixe l'\'el\'ement $(0,0)$ de $(\ZZ/n)^2$ ou, ce qui revient au m\^eme, le morphisme $\varphi \colon \ZZ/n \to (\ZZ/n)^2: x \mapsto (x,0)$. Raisonnant comme en 4.2, on obtient:

\begin{construction}[4.8]
$\mll_{00}(n)[1/n]$ s'identifie au champ alg\'ebrique classifiant les courbes elliptiques $C/S$, avec $n$ inversible sur $S$, munies d'un plongement $\varphi \colon \ZZ/n \hookrightarrow C_n$, i.e.\ d'une section $\varphi(1)$ partout d'ordre exactement $n$.
\end{construction}

\subsubsection{}
$\Gamma'_{00}(n)$ est l'image r\'eciproque du sous-groupe $H$ de $\Gl(2, \ZZ/n)$ qui fixe $(0,0)$ et la classe de $(0,1)$ dans $(\ZZ/n)^2/\ZZ/n \times \{0\}$. Des lors, $\mll'_{00}(n)[1/n]$ s'identifie au champ alg\'ebrique classifiant les courbes elliptiques $C/S$, avec $n$ inversible sur $S$, munies de $\varphi \colon \ZZ/n \hookrightarrow C_n$ et $\psi \colon \ZZ/n \to C_n/\im(\varphi)$. Le produit ext\'erieur d\'efinit un isomorphisme
\[
C_n/\im(\varphi) \xrightarrow{\sim} C_n \otimes C_n \otimes (C_n/\im(\varphi)) \xrightarrow{e_n} \mu_n.
\]
Via cet isomorphisme, se donner $\psi$ revient donc \`a se donner un isomorphisme de $\ZZ/n$ avec $\mu_n$, i.e.\ une racine primitive $n$-i\`eme de l'unit\'e sur $S$. Des lors:

\begin{construction}[4.10]
\[
\mll'_{00}(n)[1/n] \simeq \mll_{00}(n)[1/n] \times_{\Spec(\ZZ[1/n])} \Spec(\ZZ[\zeta_n, 1/n]).
\]
\end{construction}

Cet isomorphisme est le produit fibr\'e, sur $\Spec(\ZZ)$, du morphisme d'oubli: $\mll'_{00}(n) \to \mll_{00}(n)$ (d\'efini parce que $\Gamma'_{00}(n) \subset \Gamma_{00}(n)$) et du morphisme $d$ (3.20.3).

\begin{definition}[4.11]
$\mll'_0(n)[1/n]$ est le champ alg\'ebrique classifiant les courbes elliptiques g\'en\'eralis\'ees $C/S$ avec $n$ inversible sur $S$, munie d'un sous-sch\'ema en groupes $A$ localement isomorphe \`a $\ZZ/n$, et rencontrant chaque composante irr\'eductible de chaque fibre g\'eom\'etrique de $C$.
\end{definition}

Si $C$ un $m$-gone est classifi\'e par $\mll'_0(n)$, on a $n \mid m$. L'application ``oubli de $A$'' envoie donc $\mll'_0(n)[1/n]$ dans $\mll^{h*}_1$. Elle est \'etale de sorte que $\mll'_0(n)[1/n]$ est un champ alg\'ebrique lisse.

\subsubsection{}
Comme en 3.5, on d\'efinit un morphisme ``oubli de $A$ et contraction (1.3)''
\[
c \colon \mll'_0(n)[1/n] \longrightarrow \mll_1[1/n] \colon (C,A) \mapsto c(C).
\]
Pour $k$ un corps alg\'ebriquement clos et $C$ un objet de $\mll'_0(n)[1/n](k)$, l'application $\Aut(C) \to \Aut(c(C))$ est injective: pour $C$ lisse, c'est trivial; pour $C$ un $m$-gone, muni de $A \subset C_n$, on v\'erifie sur II.1.10 que $x \mapsto -x$ est le seul automorphisme non trivial de $(C,A)$. On v\'erifie comme en 3.4 que $\mll'_0(n)[1/n]$ est propre sur $\ZZ[1/n]$. Comme en 3.5, on en d\'eduit le r\'esultat suivant:

\begin{construction}[4.13]
On a $\mll_0(n)[1/n] = \mll'_0(n)[1/n]$.
\end{construction}

Des arguments analogues fournissent la

\begin{construction}[4.14]
$\mll_{00}(n)[1/n]$ est le champ alg\'ebrique classifiant les courbes g\'en\'eralis\'ees $C/S$ avec $n$ inversible sur $S$, munies d'un plongement $\delta \colon \ZZ/n \hookrightarrow C_n$, dont l'image rencontre chaque composante irr\'eductible de chaque fibre g\'eom\'etrique de $C$.
\end{construction}

Puisque $\ZZ[\zeta_n][1/n]$ est \'etale sur $\ZZ[1/n]$, on d\'eduit par ailleurs de 4.10 par normalisation la

\begin{construction}[4.15]
\[
\mll'_{00}(n) \simeq \mll_{00}(n) \times_{\Spec(\ZZ)} \Spec(\ZZ[\zeta_n]).
\]
\end{construction}

\subsubsection{}
Soit $A$ le sous-groupe de $\Gl(2, \ZZ/n)$ form\'e des homoth\'eties. Nous appellerons structures projectives de niveau $n$ les structures de niveau $A$. Elles jouent un r\^ole important dans les travaux d'Ihara. Plus g\'en\'eralement:

\begin{definition}[4.17]
Soit $C/S$ une courbe elliptique g\'en\'eralis\'ee sur $S$, avec $n$ inversible sur $S$. On suppose que les fibres g\'eom\'etriques singuli\`eres de $C$ sont des $n$-gones. Une structure projective de niveau $n$ sur $C$ est une section de $\Isom(C_n, (\ZZ/n)^2)/A$.
\end{definition}

\begin{lemma}[4.18]
Soit $(C, \eta)$ une courbe elliptique g\'en\'eralis\'ee sur un corps alg\'ebriquement clos, munie d'une structure projective de niveau $n$. Si $n \geq 2$, le seul automorphisme non trivial de $C$ qui respecte $\eta$ est $x \mapsto -x$.
\end{lemma}

Si $C$ est un $n$-gone, le lemme r\'esulte facilement de II.1.10, et vaut m\^eme si $n=1$. Si $C$ est lisse, c'est un cas particulier du lemme suivant.

\begin{lemma}[4.19]
Soient $n \geq 2$ et $g$ un automorphisme d'ordre fini d'une vari\'et\'e ab\'elienne $X$ sur un corps alg\'ebriquement clos $k$. Si $g$ induit une homoth\'etie sur $X_n$, alors $g^2 = 1$.
\end{lemma}

Par hypoth\`ese, il existe $r \in \ZZ/n$ tel que $g-r$ soit nul sur $X_n$: il existe donc un endomorphisme $h$ de $X$ tel que $g - r = n \cdot h$. Si $\zeta$ est une racine du polyn\^ome caract\'eristique de $g$, $\zeta$ est une racine de l'unit\'e et l'entier alg\'ebrique $\zeta - r$ est divisible par $n$. Si $R$ est l'anneau des entiers du corps $\QQ(\zeta)$, l'image de $\zeta$ dans $R/nR$ est donc \'egale \`a celle de $r$. Puisque $R = \ZZ[\zeta]$, c'est absurde si $\zeta \neq \pm 1$. Puisque $g$ est semi-simple, 4.19 en r\'esulte.

Du lemme 4.18 pour les $n$-gones, et d'arguments maintenant familiers, on d\'eduit le r\'esultat suivant.

\begin{construction}[4.20]
Pour $A$ comme en 4.16, $\mll_A[1/n]$ est le champ alg\'ebrique classifiant les courbes elliptiques g\'en\'eralis\'ees munies d'une structure projective de niveau $n$.
\end{construction}

Pour $(C, \eta)$ classifi\'e par $\mll_A[1/n]$, et $n \geq 2$, il r\'esulte de 4.18 que $\Aut(C, \eta)$ est constant. On en d\'eduit la proposition suivante (I.1.8).

\begin{proposition}[4.21]
Pour $A$ comme en 4.16, et $n \geq 2$, l'application
\[
\mll_A[1/n] \longrightarrow \CM_A[1/n]
\]
est \'etale.
\end{proposition}

\begin{remark}[4.22]
Pour $n=2$, $\mll_A \cap \mll_1$ est autre que $\mll_1$, et l'application
\[
\CM_A[1/2] \longrightarrow \CM_1[1/2]
\]
est fini \'etale, de degr\'e 3.
\end{remark}


\subsection{Th\'eorie transcendante (rappels)}

\subsubsection{}
Soit $E$ une courbe elliptique sur $\CC$. L'application exponentielle d\'efinit une suite exacte de groupes analytiques
\[
0 \longrightarrow H_1(E, \ZZ) \longrightarrow \Lie(E) \xrightarrow{\exp} E(\CC) \longrightarrow 0.
\]
Si $\delta \colon \ZZ^2 \xrightarrow{\sim} H_1(E, \ZZ)$ est une base de $E$, nous poserons
\[
z(\delta) = z_1/z_2 \in \CC \cup \{\infty\}
\]
(pour $L_1$ et $L_2$ deux \'el\'ements de $\Lie(E)$, avec $L_2 \neq 0$, on note $z_1/z_2$ le nombre complexe $\lambda$ tel que $L_1 = \lambda L_2$).

Si $\gamma = \begin{pmatrix} a & b \\ c & d \end{pmatrix} \in \Gl(2, \ZZ)$, on a
\[
z(\gamma\delta) = \frac{a z(\delta) + b}{c z(\delta) + d}.
\]

\subsubsection{}
L'application $\exp(n^{-1}L(x))$ induit un isomorphisme
\begin{equation}\label{eq:5.2.1}
E_n \xrightarrow{\sim} n^{-1}H_1(E, \ZZ)/H_1(E, \ZZ) \simeq (\ZZ/n)^2.
\end{equation}
Posons $X^+ = \CC - \RR$, et soit $X$ le demi-plan sup\'erieur. Si $\alpha \colon E_n \xrightarrow{\sim} (\ZZ/n)^2$ est une structure de niveau $n$, la classe du couple $(z(\delta), \alpha \circ \delta): z(\delta) \in X^+$, $\alpha \circ \delta \in \Gl(2, \ZZ/n)$, dans
\[
X^+ \times \Gl(2, \ZZ/n) / \Gl(2, \ZZ)
\]
ne d\'epend que de $(E, \alpha)$.

\subsubsection{}
Soient $H$ et $K$ comme en 3.6. La construction 5.2 d\'efinit un isomorphisme
\[
\mll_H(\CC) \simeq X^+ \times (H \backslash \Gl(2, \ZZ/n)) / \Gl(2, \ZZ).
\]
On a d'ailleurs des isomorphismes
\begin{align*}
X \times (H \backslash \Gl(2, \ZZ/n)) / \SL(2, \ZZ) & \simeq X^+ \times (H \backslash \Gl(2, \ZZ/n)) / \Gl(2, \ZZ) \\
& \simeq K \backslash X^+ \times \Gl(2, \Af_f) / \Gl(2, \hat{\ZZ}).
\end{align*}

\subsubsection{}
On sait que le morphisme de r\'eduction $\SL(2, \ZZ) \to \SL(2, \ZZ/n)$ est surjectif. Il en r\'esulte que les composantes connexes de $\mll_H(\CC)$ s'identifient aux fibres de l'application
\begin{equation}\label{eq:5.4.0}
X \times (H \backslash \Gl(2, \ZZ/n)) / \SL(2, \ZZ) \longrightarrow \det H \backslash (\ZZ/n)^*.
\end{equation}

Soit $\mu_n^*(\CC)$ l'ensemble des racines primitives $n$-i\`emes de l'unit\'e dans $\CC$. Par extension des scalaires de $\ZZ$ \`a $\CC$, (3.20.4) d\'efinit une application
\begin{equation}\label{eq:5.4.1}
d \colon \mll_H(\CC) \longrightarrow \mu_n^*(\CC)/(\det H).
\end{equation}
Via la bijection $x \mapsto \exp(2\pi ix/n): \ZZ/n \xrightarrow{\sim} \mu_n(\CC)$, (5.4.1) s'identifie \`a la restriction de cette application \`a $\mll_H(\CC)$. Les fibres g\'eom\'etriques de caract\'eristiques $0$ de (3.20.3) (3.20.4) sont donc irr\'eductibles. Le corollaire suivant r\'esulte d\`es lors du th\'eor\`eme de connexion de Zariski.

\begin{corollary}[5.5]
En toute caract\'eristique, les fibres g\'eom\'etriques de (3.20.4)
\[
\CM_H \longrightarrow \Spec(\ZZ[\zeta_n]^{\det H})
\]
sont connexes.
\end{corollary}

\begin{corollary}[5.6]
En caract\'eristique $p$ premi\`ere \`a $n$, les fibres g\'eom\'etriques de (3.20.4) sont irr\'eductibles.
\end{corollary}

Elles sont connexes et g\'eom\'etriquement unibranches (3.10 (iv)).


\subsection{Structures de niveau $H$: \'etude \`a l'infini}

Ce \S\ ne servira pas dans la suite de cet article. Nous recommandons au lecteur de l'omettre en premi\`ere lecture.

Soit $H$ un sous-groupe de $\Gl(2, \ZZ/n)$. Nous nous proposons de donner une interpr\'etation modulaire de $\mll_H[1/n]$ qui g\'en\'eralise 3.5, 4.13, 4.14, 4.15 et 4.20.

Nous d\'efinirons tout d'abord par voie modulaire un champ $\mll'_H[1/n]$, propre et lisse sur $\ZZ[1/n]$, et dont $\mll_H[1/n]$ soit un sous-champ ouvert dense. Nous d\'efinirons comme en 3.5 un morphisme fini $c \colon \mll'_H[1/n] \to \mll_1[1/n]$. Nous prouverons ensuite l'analogue de (3.5.1). La d\'emonstration de 3.5 donne alors $\mll'_H[1/n] = \mll_H[1/n]$, c'est l'interpr\'etation modulaire cherch\'ee.

\subsubsection{}
Soient $A$ un groupe isomorphe \`a $(\ZZ/n)^2$, $A' \subset A$ un sous-groupe isomorphe \`a $\ZZ/n$ et $B$ un sous-groupe contenant $A'$. Le syst\`eme $(A, A', B)$ est donc isomorphe \`a un syst\`eme $((\ZZ/n)^2, \ZZ/n \times \{0\}, \ZZ/n \times m\ZZ/n^2)$, pour $m \mid n$ convenable. Soit de plus $\mu$ un groupe isomorphe \`a $\ZZ/n$ et $\varphi_0 \colon A \to \mu$. Nous noterons $\varphi_1$ la restriction de $\varphi_0$ \`a $A' \cap B$.

L'exemple essentiel d'une telle situation est le suivant. Soient $C$ une courbe elliptique g\'en\'eralis\'ee sur un trait complet $S$, de points $\eta$ et $s$; $\bar{\eta}$ le spectre d'une cl\^oture alg\'ebrique de $k(\eta)$ et $\bar{S}$ le spectre de la cl\^oture alg\'ebrique correspondante de $k(s)$. Si $n$ est inversible sur $S$, que $C_{\eta}$ est lisse et que $\bar{C}_{\bar{s}}$ est un polygone de N\'eron \`a $mn$ c\^ot\'es, on prend
\[
\mu = \mu_n(\bar{\eta}) \simeq \mu_n(\bar{s}), \quad A = C_n(\bar{\eta}), \quad A' = C_n(\bar{s}), \quad B = (C \times_S \bar{S})^{\reg}_n
\]
et $\varphi_0 = e_n$. La restriction de $\varphi_0$ \`a $A' \cap B$ ne d\'epend que de la fibre sp\'eciale $C_{\bar{s}}$ (cf.~3.21).

\subsubsection{}
Avec les notations de 6.1, soient $I$ l'ensemble des isomorphismes $\alpha \colon A \xrightarrow{\sim} (\ZZ/n)^2$ et $\bar{I}$ l'ensemble des couples $(\alpha, \beta)$ form\'es de $\alpha \colon B \hookrightarrow (\ZZ/n)^2$ et de $\beta \colon \mu \xrightarrow{\sim} \ZZ/n$ tels que le diagramme suivant soit commutatif
\[
\xymatrix{
A' \cap B \ar@{^(->}[r] \ar[d]_{\varphi_1} & A \ar[r]^{\alpha} & (\ZZ/n)^2 \ar[d]^{\det} \\
\mu \ar[rr]^{\beta} & & \ZZ/n
}
\]
On d\'efinit $c \colon I \to \bar{I}$ par $\alpha \mapsto (\alpha|_B, \det(\alpha) \circ \varphi_0^{-1})$. Le groupe $\Gl(2, \ZZ/n)$ agit sur $I$ et $\bar{I}$ \`a gauche par $g.\alpha = g \circ \alpha$ et $g.(\alpha, \beta) = (g \circ \alpha, \det(g) \cdot \beta)$. Soit $U$ le sous-groupe de $\Aut(A)$ qui agit trivialement sur $A'$ et $A/A'$. Il agit \`a droite sur $I$ et $\bar{I}$ par $\alpha.u = \alpha \circ u$ et $(\alpha, \beta)u = (\alpha \circ (u|_B), \beta)$. Ces actions sont compatibles \`a $c$.

\begin{lemma}[6.2.2]
L'application $c$ est surjective.
\end{lemma}

La v\'erification est laiss\'ee au lecteur.

Les sous-groupes $U'$ de $U$ et les sous-groupes $B'$ de $A$ contenant $A'$ se correspondent bijectivement par
\begin{align*}
U' & \mapsto \text{sous-groupe } A^{U'} \text{ des invariants de } U'; \\
B' & \mapsto \text{sous-groupe de } U(B') \text{ agissant trivialement sur } B'.
\end{align*}
On a $\bar{I} = I/U(B)$.

Soit $\bar{\alpha} \in H \backslash \bar{I}$, image de $\alpha \in I$. L'ensemble des $\alpha' \in I$ d'image $\bar{\alpha}$ est $H\alpha U(B)$. Nous dirons que $B$ est ad\'equat si $U(B)$ est le sous-groupe de $U$ qui stabilise l'image $\bar{\alpha}$ de $\alpha$ dans $H \backslash \bar{I}$.
\[
\xymatrix{
I \ar[r]^{c} \ar[d] & \bar{I} = I/U(B) \\
H \backslash I \ar[r] & H \backslash \bar{I} = H \backslash I/U(B)
}
\]
Ceci signifie que
\begin{equation}\label{eq:6.2.3}
\alpha \text{ est le seul \'el\'ement de } H \backslash I \text{ d'image } \bar{\alpha};
\end{equation}
\begin{equation}\label{eq:6.2.4}
\text{tout } u \in U \text{ tel que } \alpha u = \alpha \text{ est dans } U(B).
\end{equation}

Notons $H \backslash I / B$ l'ensemble des $\bar{\alpha} \in H \backslash \bar{I}$ pour lesquels $B$ est ad\'equat. Il r\'esulte de la discussion pr\'ec\'edente que

\begin{lemma}[6.2.5]
On a $H \backslash I \simeq \bigsqcup_{A' \subset B' \subset A} H \backslash I / B'$ (somme sur $B'$).
\end{lemma}

\subsubsection{}
Soient donn\'es seulement un groupe $B$, un sous-groupe $A'$, avec $(A', B)$ isomorphe \`a $(\ZZ/n \times \{0\}, \ZZ/n \times m\ZZ/n^2)$ pour $m \mid n$ convenable, $\mu$ isomorphe \`a $\ZZ/n$ et $\varphi_1 \colon A' \cap B \to \mu$. Notons encore $\bar{I}$ l'ensemble des couples $(\alpha, \beta)$ form\'es de $\alpha \colon B \hookrightarrow (\ZZ/n)^2$ et de $\beta \colon \mu \xrightarrow{\sim} \ZZ/n$ tels que le diagramme 6.2.1 soit commutatif.

Le groupe $\Gl(2, \ZZ/n)$ agit \`a gauche sur $\bar{I}$ par $g.(\alpha, \beta) = (g \circ \alpha, \det(g) \cdot \beta)$ et le groupe $U$ des automorphismes de $B$ induisant l'identit\'e sur $A'$ et $B/A'$ agit \`a droite par $(\alpha, \beta)u = (\alpha \circ u, \beta)$. Nous noterons $H \backslash I / B$ l'ensemble des $\bar{\alpha} \in H \backslash \bar{I}$ tels que
\begin{enumerate}[label=\alph*)]
\item si $\bar{\alpha} \in \bar{I}$ est d'image $\bar{\alpha}$, et que $g\bar{\alpha} = \bar{\alpha}$, on a $g \in H$;
\item tout $u \in U$ tel que $\bar{\alpha} u = \bar{\alpha}$ est l'identit\'e.
\end{enumerate}
Cette notation co\"incide avec celle de 3.11 lorsque cette derni\`ere a un sens.

\subsubsection{}
Soient $S$ un sch\'ema sur lequel $n$ est inversible, et $C$ une courbe elliptique g\'en\'eralis\'ee sur $S$. On suppose que les fibres g\'eom\'etriques de $C$ sont lisses ou des polygones de N\'eron \`a $m$ c\^ot\'es, avec $n \mid m$. On se propose de d\'efinir un faisceau \'etale $\CF$ sur $S$ de formation compatible \`a tout changement de base, intuitivement d\'ecrit par les deux propri\'et\'es suivantes.
\begin{enumerate}[label=\alph*)]
\item Si $C$ est lisse sur $S$, $\CF(S)$ est l'ensemble de structures de niveau $H$ de $C$ sur $S$.
\item Supposons $S$ spectre d'un corps alg\'ebriquement clos, et que $C$ soit un polygone de N\'eron \`a $m$ c\^ot\'es ($n \mid m$). Appliquons 6.3 pour $B = C_n$, $A' = C_n^{\circ}$ et $\varphi_1 = e_n$ (cf.~6.1). On veut avoir $\CF(S) = H \backslash I / B$.
\end{enumerate}

\subsubsection{}
Soit $\CF_1(S)$ l'ensemble des syst\`emes $(B, \alpha, \beta)$ consistant en
\begin{enumerate}[label=\alph*)]
\item un sous-sch\'ema en groupes fini \'etale $B \subset C_n$;
\item $\alpha \colon B \hookrightarrow (\ZZ/n)^2$ et $\beta \colon \mu_n \xrightarrow{\sim} \ZZ/n$;
\end{enumerate}
tels que
\begin{enumerate}[label=\alph*), resume]
\item localement sur $S$, $B$ contient un sous-groupe isomorphe \`a $\ZZ/n$;
\item le diagramme
\[
\xymatrix{
A' \cap B \ar@{^(->}[r] \ar[d] & B \ar[r]^{\alpha} & (\ZZ/n)^2 \ar[d]^{\det} \\
\mu_n \ar[rr]^{\beta} & & \ZZ/n
}
\]
est commutatif;
\item pour tout point g\'eom\'etrique $\bar{s}$ de $S$, si $C_{\bar{s}}$ est singulier, $(\alpha, \beta) \bmod H$ est du type d\'ecrit en 6.4. Dans tous les cas, on demande que si $g \in \Gl(2, \ZZ/n)$ est tel que $g\alpha = \alpha$, alors $g \in H$.
\end{enumerate}

Si $C/S$ est lisse, un syst\`eme $(B, \alpha, \beta)$ d\'efinit une et une seule $H$-structure sur $C$: celle localement repr\'esent\'ee par des $\alpha \colon C_n \xrightarrow{\sim} (\ZZ/n)^2$ qui prolongent $\alpha$ et de d\'eterminant $\beta$ (l'existence locale de $\alpha$ r\'esulte de 6.2.1, l'unicit\'e mod $H$ de la condition e)).

On dira que $a$ et $b$ dans $\CF_1(S)$ sont \'equivalents si:
\begin{enumerate}[label=\arabic*)]
\item Soit $U$ l'ouvert de $S$ au-dessus duquel $C$ est lisse; $a$ et $b$ doivent d\'efinir la m\^eme structure de niveau $H$ sur $C_U/U$.
\item Soit $\bar{s}$ un point g\'eom\'etrique de $S$ tel que $C_{\bar{s}}$ soit singulier. $a_{\bar{s}}$ et $b_{\bar{s}}$ doivent \^etre conjugu\'es sous $H$. Cette condition implique que $a$ et $b$ sont conjugu\'es sous $H$ dans un voisinage \'etale $V$ de $\bar{s}$, donc que 1) est v\'erifi\'e dans $V$.
\end{enumerate}

Le faisceau $\CF$ des structures de niveau $H$ sur $C/S$ est le faisceau engendr\'e par le pr\'efaisceau $\CF_0$, quotient de $\CF_1$ par la relation d'\'equivalence pr\'ec\'edente.

On a fait ce qu'il fallait pour que $\CF$ soit de formation compatible \`a tout changement de base, admette les valeurs 6.4 et soit repr\'esent\'e par un $S$-sch\'ema $S'$ \'etale sur $S$.

\begin{definition}[6.6] \ \begin{enumerate}[label=(\roman*)]
\item Une structure de niveau $H$ sur une courbe elliptique g\'en\'eralis\'ee $C/S$ comme en 6.4 est un \'el\'ement de $\CF(S)$.
\item Le champ alg\'ebrique $\mll^{*}_H[1/n]$ sur $\ZZ[1/n]$ est le champ alg\'ebrique qui classifie les courbes elliptiques g\'en\'eralis\'ees $C/S$, ($n$ inversible sur $S$) \`a fibres g\'eom\'etriques de $C$ lisse ou des polygones de N\'eron \`a $mn$ c\^ot\'es, munies d'une structure de niveau $H$.
\end{enumerate}
\end{definition}

Par oubli de la structure de niveau $H$ on d\'efinit un morphisme \'etale $\mll^*_H[1/n] \to \mll^*_1$; $\mll^*_H[1/n]$ est donc lisse sur $\ZZ[1/n]$ et $\mll^{h*}_H[1/n]$ \'etale sur $\ZZ[1/n]$.

Par contraction (1.3 pour $n=1$) et oubli de la structure de niveau, on d\'efinit $c \colon \mll^*_H[1/n] \to \mll_1$.

\begin{theorem}[6.7] \ \begin{enumerate}[label=(\roman*)]
\item $\mll^*_H[1/n]$ est propre et lisse sur $\ZZ[1/n]$; l'image $\mll^{\infty*}_H[1/n]$ du sous-sch\'ema de non lissit\'e de la courbe universelle est fini \'etale sur $\ZZ[1/n]$.
\item Le morphisme $c$ est fini et repr\'esentable. Le champ $\mll^*_H[1/n]$ co\"incide donc (cf.~3.5) avec $\mll_H[1/n]$.
\end{enumerate}
\end{theorem}

Il reste \`a prouver que $\mll^*_H[1/n]$ est propre sur $\ZZ[1/n]$, et que $c$ est repr\'esentable, i.e.\ que, pour tout objet $C$ du champ $\mll^*_H[1/n]$ sur $k$ alg\'ebriquement clos,
\begin{equation}\label{eq:6.7}
\Aut(C) \longrightarrow \Aut(c(C))
\end{equation}
est injectif (cf.~3.5).

La propret\'e r\'esultera du crit\`ere valuatif de propret\'e et de 6.2.5. Soient $S$ un trait complet \`a corps r\'esiduel alg\'ebriquement clos, $C_{\eta}$ une courbe elliptique sur le point g\'en\'erique $\eta$ de $S$, et $\alpha$ une structure de niveau $H$ sur $C_{\eta}$. On peut supposer que le mod\`ele minimal $C'$ de $C_{\eta}$ est lisse ou un polygone de N\'eron \`a $k$ c\^ot\'es, avec $n \mid k$. Si $C'$ est lisse, la structure de niveau $H$ de $C_{\eta}$ se prolonge de fa\c{c}on unique \`a $C'$ et on pose $\bar{C} = C'$. Sinon, soit $\bar{C}''$ la courbe d\'eduire de $C'$ par contraction, comme en 1.3 ($\bar{C}''_{\bar{s}}$ \`a $n$ c\^ot\'es). Avec les notations de 6.1, soit $U(B')$ le plus grand sous-groupe de $U$ qui fixe $\bar{\alpha}$. $B'$ correspond \`a un sous-sch\'ema en groupes fini et \'etale de $\bar{C}''_n$, d'ordre $n \cdot m$, et on note $\bar{C}$ la courbe d\'eduire de $\bar{C}''$ par 1.3 (avec $n$ de loc.\ cit.\ $= m$). $\alpha$ se prolonge de fa\c{c}on unique en une structure de niveau $H$ sur $\bar{C}$, \`a l'aide de 6.2.5, et comme en 1.6. On voit que $(\bar{C}, \text{prolongement de } \alpha)$ est l'unique prolongement de $(C_{\eta}, \alpha)$ sur $S$.

Prouvons (6.7) injectif. Pour $C$ lisse, c'est clair. Sinon, le noyau de (6.7.1) est l'ensemble $A$ des automorphismes de $C$, agissant trivialement sur $\Pic^0(C)$, et respectant la structure de niveau. La condition e) dans la d\'efinition 6.5 des structures de niveau assure que $A = \emptyset$.
\qed


\section{Chapitre V. R\'eduction modulo $p$}

Dans ce chapitre, nous \'etudions la r\'eduction modulo $p$ des champs $\mll_K$ pour $K = \Gamma_0(p)$, $\Gamma_{00}(p)$ et $\Gamma(p)$. La comparaison des r\'esultats obtenus pour $\Gamma_0(p)$ et $\Gamma_{00}(p)$ nous permet de montrer que, conform\'ement \`a une conjecture de Shimura, certains sch\'emas ab\'eliens \'etudi\'es par Shimura et Casselman [5] ont bonne r\'eduction sur l'anneau des entiers d'extensions cyclotomiques convenables de $\QQ$.

\subsection{\'Etude de $\mll_0(p)$}

Nous nous proposons dans ce \S\ de donner une interpr\'etation modulaire du champ $\mll_0(p)$ lorsque $p$ est premier. Quelques r\'esultats pr\'eliminaires serviront aux paragraphes suivants.

\subsubsection{}
Soit $\CU_n$ le champ classifiant les courbes elliptiques g\'en\'eralis\'ees \`a fibres g\'eom\'etriques irr\'eductibles $C/S$, munies d'un sous-sch\'ema en groupes $A$ localement isomorphe \`a $\ZZ/n$ (pour la topologie \'etale), i.e.\ fini plat localement libre de rang $n$ de dual de Cartier \'etale.

\begin{proposition}[1.2]
Le morphisme ``oubli de $A$'' $\CU_n \to \mll_1$ est \'etale; en particulier, $\CU_n$ est un champ alg\'ebrique lisse sur $\ZZ$.
\end{proposition}

Soit $C/S$ classifi\'e par $\mll_1$. Il faut v\'erifier que le foncteur $F \colon (\Sch/S) \to (\ens)$, qui \`a $T/S$ associe l'ensemble des sous-groupes $A$ de $C_T$ localement isomorphes \`a $\ZZ/n$, est repr\'esentable par un sch\'ema $M$ \'etale sur $S$. La repr\'esentabilit\'e r\'esulte de la th\'eorie du sch\'ema de Hilbert; $M$ est \'etale car formellement \'etale (SGA 3 IX 3.6).

\subsubsection{}
Soit $\CV_n$ le champ classifiant les courbes elliptiques g\'en\'eralis\'ees $C/S$, \`a fibres g\'eom\'etriques lisses ou des $n$-gones, munies d'un sous-sch\'ema en groupes $B$ localement isomorphe \`a $\ZZ/n$ (pour la topologie \'etale), qui rencontre chaque composante irr\'eductible de chaque fibre g\'eom\'etrique.

Soit $(C,B)$ comme plus haut. Le sch\'ema en groupes $B$ agit librement sur $C$ et $C/B$ est \`a fibres g\'eom\'etriques irr\'eductibles. Dans la suite exacte
\[
0 \longrightarrow B \longrightarrow C_n \longrightarrow C_n/B \longrightarrow 0,
\]
les termes extr\^emes sont en dualit\'e de Cartier. Le couple $(C/B, C_n/B)$ est donc classifi\'e par $\CU_n$ et ceci d\'efinit
\begin{equation}\label{eq:1.3.1}
w \colon \CV_n \longrightarrow \CU_n.
\end{equation}

\begin{proposition}[1.4]
Le morphisme 1.3.1 est un isomorphisme. En particulier, $\CV_n$ est un champ alg\'ebrique lisse sur $\ZZ$.
\end{proposition}

Notons encore $w$ le morphisme de $\CV_n$ dans $\CU_n$. Une construction identique \`a celle de IV.4.4 montre que ce morphisme est un inverse de la restriction $w \colon \CV_n \to \CU_n$.

Soit $(C,A)$ sur $S$, classifi\'e par $\CU_n$. Il faut prouver qu'il existe $(\bar{C}, \bar{A})$ sur $S$, classifi\'e par $\CV_n$, un isomorphisme $\xi \colon w(\bar{C}, \bar{A}) \simeq (C,A)$, et que $(\bar{C}, \bar{A}, \xi)$ est unique \`a isomorphisme unique pr\'es. Par abus de notation, on note encore $\xi$ le morphisme compos\'e $\bar{C} \xrightarrow{\xi} C \times_{C/\bar{A}} \bar{C}/\bar{A} \to C$. La courbe \`a section marqu\'ee $(\bar{C}, e)$ sera n\'ecessairement un rev\^etement \'etale de degr\'e $n$ de $(C,e)$, muni de sa structure de courbe elliptique g\'en\'eralis\'ee.

Rappelons le lemme suivant, qui r\'esulte de SGA 4 XII.5.9 bis.

\begin{lemma}[1.5]
Soit $f \colon X \to S$ un morphisme propre \`a fibres g\'eom\'etriques connexes et $e$ une section de $f$. Si $p \colon (Y,e) \to (X,e)$ est un rev\^etement fini \'etale \`a section marqu\'ee de $(X,e)$ et que $Y/S$ est \`a fibres g\'eom\'etriques connexes, le triple $(Y,p,e)$ n'a pas d'automorphisme non trivial. Le foncteur $\CF$ sur $(\Sch/S)$ des classes d'isomorphie des triples $(Y,p,e)$ comme plus haut est repr\'esentable par un faisceau \'etale.
\end{lemma}

Appliquons ce lemme \`a $C/S$. Soit $\CJ$ le champ des $(\bar{C}, \bar{A}, \xi)$ du type voulu. Le lemme montre que les objets de $\CJ$ n'ont pas d'automorphisme non trivial, et que le foncteur $[\CJ]$ des classes d'isomorphie d'objets de $\CJ$ est le sous-foncteur de $\CF$ form\'e des $(\bar{C}, p, e)$ tels que $p^{-1}(e)$ soit localement isomorphe \`a $\ZZ/n$ et que $\bar{A}$ soit l'image de $\bar{C}_n$. Ces conditions sont ouvertes (car $\bar{A}$ est de type multiplicatif), donc $[\CJ]$ est repr\'esent\'e par $T$ \'etale sur $S$.

L\`a o\`u $C$ est lisse, on sait que $T \simeq S$. Par ailleurs, un $1$-gone sur $k$ alg\'ebriquement clos admet un unique rev\^etement \'etale par un $n$-gone. Toutes les fibres g\'eom\'etriques de $T/S$ sont donc r\'eduites \`a un point, et $T \simeq S$. Ceci prouve 1.4.

\subsubsection{}
Soit $p$ un nombre premier.

\begin{theorem}[1.6]
$\mll_0(p)$ est le champ alg\'ebrique classifiant les courbes elliptiques g\'en\'eralis\'ees $C/S$, munies d'un sous-sch\'ema en groupes $A$ localement libre de rang $p$, qui rencontre chaque composante irr\'eductible de chaque fibre g\'eom\'etrique de $C/S$.
\end{theorem}

\noindent\textit{Preuve.} Rappelons [22] qu'un sch\'ema en groupes (commutatif) localement libre de rang $p$ est tu\'e par $p$. Pour $A$ comme plus haut, on a donc $A \subset C_p$.

Notons provisoirement $\mll'_0(p)$ le champ qui \`a chaque sch\'ema $S$ associe la cat\'egorie ayant pour objets les $(C,A)$ comme en 1.6 et pour morphismes les isomorphismes. En IV.4.3, nous avons construit un isomorphisme entre $\mll_0(p)[1/p]$ et $\mll'_0(p)[1/p]$. L'isomorphisme 1.6 prolongera cet isomorphisme; raisonnant comme en IV.3.5, on trouve que 1.6 r\'esulte de la conjonction des assertions suivantes:

\begin{enumerate}[label=\Alph*)]
\item $\mll'_0(p)$ est un champ alg\'ebrique.
\item $\mll'_0(p)$ est propre sur $\Spec(\ZZ)$.
\item Le morphisme $c \colon \mll'_0(p) \to \mll_1 \colon (C,A) \mapsto c(C)$ d\'efini par IV.1.4 est fini et repr\'esentable.
\item $\mll'_0(p)$ est normal.
\item $\mll'_0(p)$ est de Cohen-Macaulay.
\end{enumerate}

\medskip

\noindent\textbf{(1) Points \`a l'infini de $\mll'_0(p)$.}

Soit $C$ le polygone de N\'eron \`a un c\^ot\'e sur $\Spec(\ZZ)$. On a $C_p \simeq \mu_p$ et $(C, C_p)$ est une section de $\mll'_0(p)$ sur $\ZZ$.

Soit $C$ le $p$-gone sur $\Spec(\ZZ)$. On a $C^{\reg} \simeq \Gm \times \ZZ/p$, $C_p \simeq \mu_p \times \ZZ/p$ et $(C, e \times \ZZ/p)$ est une section de $\mll'_0(p)$ sur $\ZZ$.

Soit $k$ un corps alg\'ebriquement clos. On v\'erifie que les seuls objets $(C,A)$ de $\mll'_0(p)(k)$, avec $C$ singulier, sont ceux d\'ecrits ci-dessus.

\medskip

\noindent\textbf{(2) Preuve de A (repr\'esentabilit\'e de $\mll'_0(p)$).}

Soient $\CU_1$ et $\CU_2$ les deux ouverts suivants de $\mll'_0(p)$: $\CU_1$ classifie les courbes $(C,A)$ comme en 1.6, les fibres g\'eom\'etriques de $C$ \'etant irr\'eductibles, et $\CU_2$ classifie les courbes $(C,A)$, les fibres g\'eom\'etriques de $C$ \'etant lisses ou des $p$-gones.

\begin{lemma}[1.7]
Le morphisme ``oubli'' $\CU_1 \to \mll_1$ est repr\'esentable.
\end{lemma}

\noindent\textit{Preuve.} Il faut voir que pour $C/S$ une courbe elliptique g\'en\'eralis\'ee \`a fibres g\'eom\'etriques irr\'eductibles sur $S$, le foncteur suivant $F \colon (\Sch/S) \to (\ens)$ est repr\'esentable. \`A $T/S$, $F$ associe l'ensemble des sous-sch\'emas en groupes finis localement libres de rang $p$ de $C_T$. La repr\'esentabilit\'e de $F$ r\'esulte de la th\'eorie des sch\'emas de Hilbert.

D'apr\`es 1.7 et IV.1.2, $\CU_1$ est un champ alg\'ebrique. Nous v\'erifions ci-dessous que $\CU_1$ et $\CU_2$ sont isomorphes. Le champ $\CU_2$ est donc alg\'ebrique, ainsi que $\mll'_0(p)$, r\'eunion de $\CU_1$ et $\CU_2$ d'apr\`es (1).

\subsubsection{}
Soit $(C,A)$ classifi\'e par $\CU_1$. D'apr\`es II.1.20, $C_p$ est localement libre de rang $p^2$. On v\'erifie fibre par fibre que $A$ agit librement sur $C$, et ceci permet de d\'efinir
\begin{equation}\label{eq:1.8.0}
w(C,A) = (C/A, C_p/A).
\end{equation}

\begin{lemma}[1.8.2]
Le morphisme (1.8.1) est un isomorphisme.
\end{lemma}

Avec les notations 1.1 et 1.3, on a
\[
w \colon \CU_1 \xrightarrow{w_1} \CV_p \xrightarrow{w} \CU_p \xrightarrow{w_2} \CU_1.
\]
La construction IV.4.4 montre encore que $w$ induit une involution de $\mll'_0(p)$. En 1.4, on a vu que $w$ induit un isomorphisme $\CV_p \simeq \CU_p$, et 1.8.2 en r\'esulte.

Ceci ach\`eve la d\'emonstration de (A). La d\'emonstration nous a fourni la construction suivante.

\begin{construction}[1.9]
Par recollement de $w \colon \CU_1 \xrightarrow{\sim} \CU_1$ et de $w \colon \CU_2 \xrightarrow{\sim} \CU_2$, on obtient une involution $w$ de $\mll'_0(p) = \mll_0(p)$, qui prolonge l'involution (IV.4.4).
\end{construction}

La r\'eunion de $\CU_p$ et $\CV_p$ est l'ouvert de $\mll'_0(p)$ correspondant aux courbes non supersinguli\`eres de caract\'eristique $p$. Les propositions 1.2 et 1.4 fournissent donc:

\begin{proposition}[1.10]
$\mll'_0(p)$ est lisse sur $\Spec(\ZZ)$ en dehors des points supersinguliers de caract\'eristique $p$.
\end{proposition}

\medskip

\noindent\textbf{(3) Preuve de (B) (propret\'e de $\mll'_0(p)$).}

D'apr\`es (1.10), $\mll^{\circ}_0(p)$ est dense dans $\mll'_0(p)$. Le crit\`ere valuatif de propret\'e nous ram\`ene alors \`a d\'emontrer le lemme suivant.

\begin{lemma}[1.11]
Soit $(S, \eta, s)$ un trait complet \`a corps r\'esiduel alg\'ebriquement clos. Soit $(C,A)$ un objet de $\mll'_0(p)(\eta)$. Il existe une extension finie $k(\eta')$ de $k(\eta)$ telle que, notant $S'$ le normalis\'e de $S$ dans $k(\eta')$, l'image $(C',A')$ de $(C,A)$ dans $\mll'_0(p)(\eta')$ provienne par changement de base d'un objet $(\bar{C}, \bar{A})$ de $\mll'_0(p)(S')$, unique \`a isomorphisme unique pr\'es.
\end{lemma}

\noindent\textit{Preuve:} Quitte \`a faire une extension pr\'alable de $k(\eta)$, on peut supposer que $C$ a r\'eduction semi-stable et que la fibre sp\'eciale de son mod\`ele minimal $\tilde{C}$ soit lisse ou un $m$-gone, avec $p \nmid m$ (IV.1.6). Le groupe $C_p$ est alors fini sur $S$ (II.1.20) et l'adh\'erence sch\'ematique $\bar{A}$ de $A$ est un sous-sch\'ema en groupes fini localement libre de rang $p$ de $\tilde{C}$. On obtient $(\bar{C}, \bar{A})$ en contractant les composantes de $\tilde{C}_s$ qui ne rencontrent pas $\bar{A}$ (IV.1.3).

L'unicit\'e se d\'emontre comme (IV.1.6 (iv)): \`a $\bar{C}$ g\'en\'erale, on obtient le mod\`ele minimal $\tilde{C}$ de $C$ en \'eclatant successivement les singularit\'es, comme en (IV.1.6 (iv)). L'image r\'eciproque de $\bar{A}$ est l'adh\'erence sch\'ematique $A$ de $A$ et il est clair que $\bar{C}$ est obtenu en contractant les composantes irr\'eductibles de $\tilde{C}_s$ dont l'intersection avec $A$ est vide, d'o\`u l'unicit\'e.

\medskip

\noindent\textbf{(4) Preuve de C ($c$ fini repr\'esentable).}

Pour v\'erifier que les fibres g\'eom\'etriques de $c$ sont finies, nous distinguerons trois cas.

\begin{enumerate}[label=\alph*)]
\item \textbf{Fibres \`a l'infini:} sur $k$ alg\'ebriquement clos, il n'y a qu'un nombre fini (deux) de classes d'isomorphie de couples $(C,A)$ classifi\'es par $\mll'_0(p)$, avec $C$ singulier.

\item \textbf{Fibres \`a distance finie:} Soient $C$ une courbe elliptique sur $k$ et $T$ la fibre de $c$ en le point g\'eom\'etrique $\Spec(k) \to \mll_1$ d\'efini par $C$. $T$ repr\'esente le foncteur des sous-sch\'emas en groupes de rang $p$ de $C_p$. Distinguons les cas $C_p \simeq \ZZ/p \times \mu_p$ et $\car(k) \neq p$, $C_p \simeq \CL_p^2$.

\begin{enumerate}[label=b\arabic*)]
\item On a une suite exacte $0 \to \mu_p \to C_p \to \ZZ/p \to 0$ (1).

Pour tout sch\'ema $S$, et $A \subset C_p$ localement libre de rang $p$, on a alors $S = S' \amalg S''$, avec $A \simeq \ZZ/p$ sur $S''$ et $A \subset \mu_p$, donc $A \simeq \mu_p$ sur $S'$. On a donc $T = T' \cup T''$. $T'$ est r\'eduit, et r\'eduit au point d\'efini par $\mu_p \subset C_p$. $T''$ repr\'esente le foncteur des trivialisations $\ZZ/p \xrightarrow{\sim} C_p/\mu_p$ de la suite exacte (1). C'est un torseur sous $\mu_p$. Le sch\'ema $T$ a donc $p+1$ (resp.\ 2) points g\'eom\'etriques si $\car(k) \neq p$ (resp.\ $= p$).

\item Tout sous-groupe de $\CL_p^2$ est infinit\'esimal. Si d'ordre $p$, il est \'egal au noyau de Frobenius $\alpha_p$ et $T$ n'a qu'un point g\'eom\'etrique.
\end{enumerate}
\end{enumerate}

Puisque $\mll'_0(p)$ est propre, $c$, propre \`a fibres finies, est fini. Pour prouver $c$ repr\'esentable, on proc\'ede comme en IV.3.5: on v\'erifie que pour $(C,A)$ singulier, $\Aut(C,A) = \{1\}$, de sorte que $\Aut(C,A) \hookrightarrow \Aut(c(C))$.

\medskip

\noindent\textbf{(5) Le point clef de la d\'emonstration est le lemme suivant.}

\begin{lemma}[1.12]
Le morphisme $c$ est plat.
\end{lemma}

\noindent\textit{Preuve:} Nous prouverons 1.12 \`a l'aide du crit\`ere suivant, applicable car $\mll_1$ est r\'eduit.

\begin{lemma}[1.13]
Soit $f \colon X \to S$ un morphisme fini de sch\'emas noeth\'eriens, avec $S$ r\'eduit. Si le rang des fibres g\'eom\'etriques de $f$ est constant, alors $f$ est plat.
\end{lemma}

On peut supposer que $S$ est le spectre d'un anneau local $A$ et que $X$ est le spectre d'une $A$-alg\`ebre finie $B$. Soient $\mg$ l'id\'eal maximal de $A$ et $e_1, \ldots, e_n$ une base du $A/\mg$-espace vectoriel $B/\mg B$. Relevons les $e_i$ en des \'el\'ements $\tilde{e}_i \in B$. D'apr\`es le lemme de Nakayama, les $\tilde{e}_i$ engendrent le $A$-module $B$. Si $\sum \lambda_i \tilde{e}_i = 0$, avec $\lambda_i \in A$, il r\'esulte de l'hypoth\`ese que les $\lambda_i$ sont nuls modulo tout id\'eal premier, donc nuls puisque $A$ est r\'eduit; le $A$-module $B$ est donc libre de base les $\tilde{e}_i$.

V\'erifions que le morphisme 1.12 v\'erifie l'hypoth\`ese de 1.13.

Soit donc $C$ une courbe elliptique sur un corps alg\'ebriquement clos $k$. La fibre $T$ de $c$ en le point g\'eom\'etrique $\Spec(k) \to \mll_1$ d\'efini par $C$ repr\'esente le foncteur des sous-sch\'emas en groupes de rang $p$ de $C_p$. Prouvons qu'il est de rang $p+1$.

Si $C_p \simeq \ZZ/p \times \mu_p$, on a vu en (4) que $T$ est isomorphe \`a $\Spec(k) \amalg \mu_p$, donc de rang $p+1$.

Il reste \`a traiter le cas o\`u $k$ est de caract\'eristique $p$ et que $C_p \simeq \CL_p^2$. Dans ce cas, $T$ repr\'esente le foncteur des sous-sch\'emas localement libre de rang $p$ de $C_p$, contenus dans $\CL_p^2$.

Un sous-sch\'ema localement libre de rang $p$ de $C_p$, sur une base quelconque, est d\'efini par une \'equation
\begin{equation}\label{eq:star}
x^p + \sum_{i=0}^{p-1} a_i x^i = 0.
\end{equation}
Pour que $A$ soit stable par l'homoth\'etie $x \mapsto nx$ ($n \in (\ZZ/p)^*$), il faut et il suffit que (\ref{eq:star}) soit identique \`a l'\'equation
\[
x^p + \sum_{i=0}^{p-1} n^{p-i} a_i x^i = 0.
\]
En particulier, si $A$ est un sous-sch\'ema en groupes de $C_p$, son \'equation se r\'eduit \`a
\[
x^p - a x = 0.
\]
Pour que $A$ soit sous-sch\'ema de $\CL_p^2$, il faut et il suffit que le premier membre de cette \'equation divise $X^{p^2}$. Puisque
\[
(X^p - aX) \cdot \Bigl(\sum_{i=0}^{p-1} a^{p-1-i} X^{pi}\Bigr) = X^{p^2} - a^p X,
\]
tel est le cas si et seulement si $a^{p-1} = 0$. On a
\[
T \simeq \Spec(k[a]/(a^{p-1})) \amalg \Spec(k)
\]
et ceci ach\`eve la d\'emonstration de 1.12.

Puisque $c$ est fini et plat et $\mll_1$ r\'egulier, $\mll'_0(p)$ est de Cohen-Macaulay (EGA IV.15.4.2), purement de dimension relative un sur $\Spec(\ZZ)$. \`A l'infini, $\mll'_0(p)$ est m\^eme lisse (1.10); $\mll'_0(p)$ tout entier est donc de Cohen-Macaulay et, d'apr\`es 1.10 encore, lisse sur $\Spec(\ZZ)$ en dehors d'un ensemble fini (de codimension 2). D\`es lors, (D) r\'esulte du crit\`ere de normalit\'e de Serre (EGA IV.5.8.6).
\qed

\subsubsection{Variantes}
Soit $n$ un entier premier \`a $p$, $H \subset \Gl(2, \ZZ/n)$ et $K$ l'image r\'eciproque de $H$ dans $\Gl(2, \hat{\ZZ})$. Le th\'eor\`eme 1.6 se g\'en\'eralise aux $\mll_0(p)[1/n]$:

(1.14.1) $\mll(n) \cap \mll_0(p)[1/n]$ classifie les courbes elliptiques g\'en\'eralis\'ees $C/S$, de fibres g\'eom\'etriques lisses, des $n$-gones ou des $np$-gones, munies d'un isomorphisme $\alpha \colon C_n \xrightarrow{\sim} (\ZZ/n)^2$ et d'un sous-groupe localement libre de rang $p$ $H$ de $C$ tel que $C_n \cdot H$ rencontre chaque composante g\'eom\'etrique de chaque fibre g\'eom\'etrique.

(1.14.2) Pour $n$ sans facteur carr\'e, $\mll_0(n)$ classifie les courbes elliptiques g\'en\'eralis\'ees $C/S$ munies d'un sous-groupe localement libre de rang $n$ $H$ de $C$ qui rencontre chaque composante g\'eom\'etrique de chaque fibre g\'eom\'etrique.

\subsubsection{}
Soit $S$ un sch\'ema de caract\'eristique $p$ et $E$ une courbe elliptique sur $S$. \`A $E$, on associe les deux objets suivants de $\mll_0(p)(S)$:
\begin{align*}
\varphi_1(E) &= (E, \Ker(F \colon E \to E^{(p)})); \\
\varphi_2(E) &= w\varphi_1(E) = (E^{(p)}, \Ker(V \colon E^{(p)} \to E)).
\end{align*}
Dans ces formules, $F$ d\'esigne le morphisme de Frobenius et $V$ son transpos\'e, la Verschiebung: $FV = VF = p$. Ces constructions d\'efinissent un diagramme
\[
\xymatrix{
& \mll_0(p) \ar[dr]^{c} \ar[dl]_{c \circ w} \\
\mll_1 & & \mll_1
}
\]
Les compos\'es $c \circ \varphi_1 \colon E \mapsto (E, \Ker F) \mapsto E$ et $c \circ w \circ \varphi_1$ sont l'identit\'e. Les compos\'es $c \circ w \circ \varphi_1 \colon E \mapsto (E^{(p)}, \Ker V) \mapsto E^{(p)}$ sont l'endomorphisme de Frobenius de $\mll_1 \otimes \FF_p$.

\end{document}
